\documentclass[11pt]{article}
\usepackage[margin=0.9in]{geometry}
\usepackage{amsmath,amssymb,amsthm,booktabs,mathrsfs}
\usepackage[T1]{fontenc}
\usepackage{lmodern,microtype}
\usepackage[colorlinks=true,linkcolor=blue,urlcolor=blue]{hyperref}
\newtheorem{theorem}{Theorem}[section]
\newtheorem{lemma}[theorem]{Lemma}
\newtheorem{proposition}[theorem]{Proposition}
\newtheorem{corollary}[theorem]{Corollary}
\theoremstyle{remark}
\newtheorem{remark}[theorem]{Remark}
\newcommand{\norm}[1]{\left\|#1\right\|}
\title{Toward a four-dimensional Yang--Mills construction\\
\large Notebook 32: determinant composition and finite coupling tuning}
\author{Research working document}
\date{17 October 2026}
\setlength{\emergencystretch}{5em}
\begin{document}
\maketitle
\begin{abstract}
We identify the actual leading quantum correction and the precise
drift input needed for finite ultraviolet tuning. All statements
use the composite ordinary constraint, its pivot coordinates,
the absolute contour penalty and the retained Gaussian of the
preceding construction.

First, the finite classical Maxwell leakage tends uniformly to
zero along a cofinal nested period family, even when the preceding
homogeneous depth varies. The actual finite quantum map is jointly
continuous in its interaction class and its positive real coupling;
the continuity argument requires no continuous choice of small
representatives.

Second, we derive an explicit convergent trace formula for the
Maxwell coefficient of the pivot and Gaussian determinants.
Global color invariance removes the second stationary response
only after the full scalar chain rule is taken. The constraint,
current, gauge and logarithmic-chart curvature terms remain in
the precision derivatives. At a fixed regulator, exact Haar
pushforward composition and finite-dimensional Laplace asymptotics
prove a cocycle for these actual leading determinants. Exact
infinite classical Maxwell preservation then makes their
coefficients additive across scales.

An exact finite gauge and coarea identity also rewrites the
leading correction as vector, based-ghost and coarse-constraint
determinants. The coarse factors remain explicit. A flat
Rayleigh test shows that the based ghost alone has no uniform
gap at the final physical scale, so a reference comparison
must retain the coupled cancellations. The coarse constraint
covariance is identified with a shifted actual effective-Hessian
inverse. For a fixed sufficiently large positive shift, its
Maxwell contribution is uniformly bounded in depth and volume.

Finally, under an explicitly unproved positive averaged drift
hypothesis for the actual finite leading coefficients, a continuous
shooting argument tunes arbitrarily long finite local histories
to one fixed positive weak coupling. The auxiliary clipping is
proved absent on the selected histories. This is a conditional
tuning theorem, not an evaluation of the drift. A uniform positive
block margin in every varying-depth structural limit would suffice:
qualitative compactness transfers that margin to the actual finite
regulators without a convergence rate.

The remaining one-loop task is a bound on the difference between
the actual composite coefficient and an explicitly normalized
reference logarithm, uniform in depth. Mixed and large-field
histories, continuum observables, nontriviality and reconstruction
remain part of the full existence argument.
\end{abstract}

The objective remains a nontrivial four-dimensional Yang--Mills
quantum theory. This notebook continues the actual homogeneous
local construction of Notebooks 15--31. Its finite local density
identities do not by themselves represent the entire Wilson
measure. Results and conditional inputs are distinguished in
the statements below.


% No preamble. Prefix fml-. Qualitative uniform finite leakage on cofinal periods.
\section{Uniform disappearance of finite classical marginal leakage}
\label{fml-section}

This section gives a qualitative uniform statement along the
standard nested compatible period family. It does not assert
exact finite marginal preservation or an exponential convergence
rate. The preceding classical depth may vary arbitrarily with
the period. The proof uses the already established exact
preservation of the true infinite Maxwell coefficient in each
joint classical system, rather than identifying different
limiting systems.

\subsection{Period family and statement}
\label{fml-setup-section}

Fix the primitive ratio \(\ell\), an admissible composite ratio
\(R=\ell^s\), auxiliary grouping \(J\), carrier exponent \(a_0\),
and one positive common analytic radius \(b\). All these data
are fixed. Let \(\mathcal P\) be a compatible nested family of
cubic output periods tending to infinity and satisfying
\begin{equation}
 \text{for each fixed }q\ge0,\quad
 R^q\mid N\quad\hbox{for all sufficiently late }N\in\mathcal P.
 \label{fml-divisibility}
\end{equation}
All grouping and source-geometry divisibilities are also
imposed. The standard nested source tori have this property;
for example \(N=2\ell^k\) does when the fixed grouping is
compatible. Arbitrary growing multiples of \(J\) with bounded
\(\ell\)-adic valuation are not covered by this argument.

Let \(\mathcal T_{M,R,N}\) be the actual pre-gauge classical
substitution from the input torus of current period \(RN\)
to the output torus of period \(N\). Use the actual finite
short-lift Maxwell ratios, each normalized by its own actual
finite homogeneous reference, and define
\begin{equation}
 \Lambda_{M,R,N}
   =\beta_{MR,N}\mathcal T_{M,R,N}-\beta_{M,RN}.
 \label{fml-leakage}
\end{equation}
This is a linear functional on the invariant finite activity
quotient at the common \(a_0,b\). On an extracted input it
is the finite classical marginal leakage.

\begin{theorem}[Qualitative uniform finite leakage]
\label{fml-theorem}
For all sufficiently large lower period thresholds \(K\), set
\begin{equation}
 \omega_{R,J,b}(K)
   =\sup_{\substack{M,\ N\in\mathcal P\\N\ge K}}
      \ \sup_{\substack{x\ {\rm invariant}\\\|x\|\le b^2}}
                     |\Lambda_{M,R,N}x|.
 \label{fml-modulus}
\end{equation}
Then
\begin{equation}
                         \omega_{R,J,b}(K)\longrightarrow0
                         \quad(K\longrightarrow\infty).
 \label{fml-vanishing}
\end{equation}
The supremum allows every permitted preceding homogeneous
depth \(M\), with no fixed-depth or bounded-depth restriction.
\end{theorem}

The function in \eqref{fml-modulus} is finite and decreasing
above the reference thresholds. Indeed the uniform finite
moment bound and fixed-\(R\) transport bound give
\begin{equation}
 \omega_{R,J,b}(K)\le C_\beta(B_R+1).
 \label{fml-crude-bound}
\end{equation}
The factor \(b^{-2}\) in the dual moment norm cancels the
input normalization \(b^2\). The theorem asserts a
qualitative improvement of this bound, not a rate.

\subsection{Compactness of varying finite input classes}
\label{fml-input-compactness}

Take any sequence \(M_n\), periods \(N_n\in\mathcal P\)
tending to infinity, and invariant finite input classes
\(x_n\) with \(\|x_n\|\le b^2\). Choose representatives
\(E_n\) on the input tori with
\begin{equation}
                         \|E_n\|_{a_0,b}\le b^2+\epsilon_n,
                         \qquad \epsilon_n\downarrow0.
 \label{fml-near-minimal}
\end{equation}
Average each family over the finite translations by
\(J\)-cells, translating its carriers and arguments
together. Its anchored norm does not increase, because
the translations permute the root sums. Its normalized
total remains exactly the same by the total translation
invariance of \(x_n\). The resulting array is covariant
under \(J\)-cell translations, while the total retains
all original unit-translation, point and local gauge
identities. No coefficientwise gauge invariance is needed.

Apply the proved analytic-interaction compactness argument
at the same open radius \(b\). On every fixed embedded
connected carrier, finite-dimensional Montel compactness
gives local uniform convergence of its coefficient function.
A countable diagonal extraction gives an infinite family
\(E_\infty\). Fatou on finite rooted carrier collections,
followed by radii increasing to \(b\), yields
\begin{equation}
                         \|E_\infty\|_{a_0,b}\le b^2.
 \label{fml-limit-norm}
\end{equation}
There is no analytic-radius loss.

Vacuum-normalized totals converge on compactly supported
fields. The anchored cardinality reserve makes the
large-carrier remainder uniformly small, and every fixed
carrier eventually embeds without wrapping. Local gauge
and spatial identities pass on their common compact-field
domains by the same total-continuity argument. Thus
\(E_\infty\) represents an invariant infinite class
\(x_\infty\).

Fixed-order total kernels converge locally, in particular
\begin{equation}
                K_{2,n}^{\rm in}\longrightarrow K_2(x_\infty).
 \label{fml-input-kernel-limit}
\end{equation}
They have the common exponential row and moment bounds
following from the activity norm. No numerical value of
\(M_n\) enters this input compactness argument.

\subsection{Joint extraction when the classical depth varies}
\label{fml-varying-depth-section}

The classical maps and references must be extracted jointly
with the input interactions. For every fixed depth shift
\(j\ge0\) and endpoint power \(L=R^q\), use the actual finite
data at depth \(M_nR^j\), on compatible periods obtained
from \(RN_n\) by finitely many divisions by \(R\) and \(L\).
Condition \eqref{fml-divisibility} ensures that all these
data exist for sufficiently large \(n\), and their periods
still tend to infinity.

For each fixed endpoint \(L\), an observed current output
fiber contains only \(L^4\) times a fixed number of link
and color coordinates. This dimension is independent of
the preceding refinement \(M_n\). The previously proved
classical inputs are uniform in that refinement:
\begin{itemize}
 \item a positive common holomorphic radius;
 \item strict support and cardinality bounds;
 \item complete \(C^{1,\alpha}\) bounds and the exact
       \(1/L\) pre-gauge amplitude;
 \item actual finite reference transport and minimum
       comparison identities;
 \item the controlled finite near-identity gauge comparisons.
\end{itemize}
Apply analytic compactness to the reference families and
complete-output-family compactness to each fixed-endpoint
map. One diagonal subsequence contains every depth shift
and every endpoint power in this countable collection.

Denote the resulting references by \(S_j^*\) and the maps
by \(\mathcal T_{j,q}^*\), with endpoint \(R^q\). The index
\(j\) records a relative depth shift. No assertion
\(S_j^*=S_{M_*R^j}\) for a finite \(M_*\) is made.
The uniform flat Maxwell normalization and exponential
kernel bounds give
\begin{equation}
                           \beta(S_j^*)=1.
 \label{fml-reference-normalization}
\end{equation}

The fixed-radius reference bounds, exact transport and
scalar semigroup pass to this joint system. For compact
fields the anchored tails justify passage of normalized
reference values, without passage of an extensive vacuum
energy. The scalar semigroup on every invariant input
class uses the actual finite small connecting gauges
and their endpoint estimates. These estimates are uniform
in \(M_n\); their compact extraction and the summable-gauge
argument therefore apply to this varying-depth sequence
as well. They never require the numerical label \(M_n\)
to be fixed.

Equivalently, for each fixed endpoint pair, pass the finite
comparison germ together with its actual controlled gauge
data, and then use holomorphic continuation on the
common field domain. The comparison germ can be smaller
for a large endpoint; the uniform pre-gauge domain remains
the same. The support estimate at fixed classical reserve
is uniform in all endpoint powers, and no high-jet
field-dilation radius is used here.

The resulting system satisfies
\begin{equation}
 \mathcal T_{j,q}^*S_j^*=S_{j+q}^*,\qquad
 \mathcal T_{j+q,r}^*\mathcal T_{j,q}^*
                         =\mathcal T_{j,q+r}^*.
 \label{fml-limit-semigroup}
\end{equation}
It also has the uniform quadratic-only spatial estimate
at fixed \(b\) for arbitrarily large endpoints \(R^q\).
It therefore satisfies precisely the hypotheses of
the exact infinite Maxwell-preservation proof, which gives
\begin{equation}
                   \beta(\mathcal T_{0,1}^*x)=\beta(x)
                   \quad\hbox{for every invariant }x.
 \label{fml-limit-preservation}
\end{equation}

This is a varying-depth structural limit, not an invocation
of a fixed-\(M\) limit theorem. Different limiting systems
need not coincide. Every possible system extracted this
way has \eqref{fml-limit-preservation}, which is the only
property needed below.

\subsection{The transported Hessian and its moment tails}
\label{fml-Hessian-section}

Let
\begin{equation}
                       H_n=Da_{M_n,R,N_n}(0)
 \label{fml-linear-response}
\end{equation}
be the actual linearized pre-gauge angle map. It takes
output current-angle directions on period \(N_n\) to
input current-angle directions on period \(RN_n\).
Its transpose in the scalar Hessian formula is the
ordinary counting-coordinate transpose: all physical
rescaling is already in this angle map.

The complete support-family estimate and Cauchy at zero
give
\begin{equation}
 |H_n(b,c)|
       \le C_R e^{-\kappa_R d(\operatorname{root}(b),c)}.
 \label{fml-response-envelope}
\end{equation}
The root is the final-unit cell containing the input
bond after division of its coordinates by \(R\).
Fixed internal link and color conventions are absorbed
into \(C_R\). One may retain the amplitude \(1/R\), but
\(R\) is fixed in this proof.

Each root contains only \(O(R^4)\) input coordinates.
Thus \eqref{fml-response-envelope} gives both row and
column exponential Schur bounds, with constants independent
of \(M_n\) and volume. Joint local holomorphic convergence
implies local convergence \(H_n\to H_*\).

The total first derivative of each input vanishes by
global semisimple color invariance, and the actual
pre-gauge map vanishes at zero. Consequently the exact
finite Hessian chain rule is
\begin{equation}
                       K_{2,n}^{\rm out}
                           =H_n^T K_{2,n}^{\rm in}H_n.
 \label{fml-Hessian-product}
\end{equation}
There is no term involving the second derivative of
the background map.

The exponential envelopes for the two response legs
and the input kernel give a uniform output exponential
row bound. To see this, connect the two output source
roots through the intermediate input roots. A fixed
fraction of the minimum response and input decay
exponent pays the output separation; the remaining
geometric sums are bounded by their infinite-lattice
sums. The \(O(R^4)\) fiber multiplicity is fixed. No
factor from \(M_n^4\) fine coordinates appears.

These bounds also prove local convergence of
\eqref{fml-Hessian-product}: truncate its two intermediate
root sums to a fixed neighborhood, pass the finite
sum entrywise, and then remove the truncation using
the common exponential majorant. Hence
\begin{equation}
 K_{2,*}^{\rm out}
       =H_*^T K_2(x_\infty)H_*,
 \label{fml-output-kernel-limit}
\end{equation}
which is the Hessian of
\(\mathcal T_{0,1}^*x_\infty\). The same estimates give
uniformly summable second displacement moments.

\subsection{Convergence of the actual finite Maxwell ratios}
\label{fml-moment-section}

For locally converging torus kernels with a uniform
exponential row bound, the symmetric short-lift second
moments converge to the true infinite second moment.
First restrict to a fixed displacement radius; that
finite sum converges locally. The remaining second
moment is bounded uniformly by an exponentially decaying
tail in that radius. The short-lift cutoff
\(\lfloor N_n/16\rfloor\) tends to infinity, so it
produces no residual term. Fixed midpoint shifts are
handled by the same majorant.

The finite reference denominators tend uniformly to one
by the established flat \(q_M\) periodization estimate,
including for varying \(M_n\). Therefore
\begin{align}
 \beta_{M_n,RN_n}(x_n)
                 &\longrightarrow\beta(x_\infty),\notag\\
 \beta_{M_nR,N_n}(\mathcal T_{M_n,R,N_n}x_n)
                 &\longrightarrow
                     \beta(\mathcal T_{0,1}^*x_\infty).
 \label{fml-ratio-limits}
\end{align}
These are the actual finite reference-normalized ratios.
Their full nonlinear references have not been replaced
by canonical periodizations, and no extensive total
action limit is used.

\subsection{Uniform contradiction and consequences}
\label{fml-proof-section}

\begin{proof}[Proof of Theorem~\ref{fml-theorem}]
If \eqref{fml-vanishing} failed, there would be
\(\delta>0\) and sequences as above with
\[
                     |\Lambda_{M_n,R,N_n}x_n|\ge\delta
\]
for every \(n\). Extract the joint subsequence constructed
in the preceding sections. By \eqref{fml-ratio-limits},
the limiting difference is
\[
 \beta(\mathcal T_{0,1}^*x_\infty)-\beta(x_\infty)=0
\]
by \eqref{fml-limit-preservation}. This contradicts the
lower bound and proves the theorem.
\end{proof}

On extracted inputs the conclusion reads
\begin{equation}
 \sup_{\substack{M,\ N\in\mathcal P,\ N\ge K\\
           \beta_{M,RN}(x)=0,\ \|x\|\le b^2}}
       |\beta_{MR,N}(\mathcal T_{M,R,N}x)|
                            \longrightarrow0.
 \label{fml-extracted-conclusion}
\end{equation}
The unextracted statement \eqref{fml-modulus} is slightly
stronger and needs no projection-loss factor. By linear
homogeneity it gives, for arbitrary invariant inputs,
\begin{equation}
 |\Lambda_{M,R,N}x|
       \le \omega_{R,J,b}(K)b^{-2}\|x\|,
       \qquad N\in\mathcal P,\ N\ge K.
 \label{fml-scaled-bound}
\end{equation}

For a prescribed per-step drift tolerance, choose \(K\)
so that \(\omega_{R,J,b}(K)\) is sufficiently small.
Every step of a permitted finite history whose output
period remains at least \(K\) then has the same leakage
bound, uniformly in the preceding depth and the length
of that history. On the established envelope
\(\|x\|\le C_{\rm int}b^2\), multiply the tolerance by
\(C_{\rm int}\), or choose \(K\) so that
\(\omega_{R,J,b}(K)\) is the desired tolerance divided
by this fixed constant.

\subsection{Qualifications}
\label{fml-scope-section}

The result is qualitative. It supplies no exponential
rate, no summability of leakage over infinitely many
changing periods, and no exact zero at finite volume.
The per-step consequence is a magnitude bound, not a
sign or an accumulated-drift assertion.

The cofinal divisibility assumption
\eqref{fml-divisibility} is part of the theorem. An
extension to arbitrary increasing compatible multiples
with bounded endpoint divisibility would require an
additional uniform period-extension or local-volume
comparison. It does not follow merely from \(N\to\infty\).

The limiting systems have not been identified with
one another. Exact infinite Maxwell preservation holds
in each separately and rules out every nonzero leakage
subsequence. No quantum drift sign, determinant value,
infinite physical running trajectory, or global
measure interpretation follows from this compactness
argument.


% No preamble. Prefix fqc-. All equation references are internal.
\section{Continuity in the real coupling at a finite regulator}
\label{fqc-section}

This section concerns the actual homogeneous local macrostep and its
specified common-domain continuation. It supplies continuity needed
for a finite endpoint argument. It proves no coupling sign, endpoint
crossing, or identification with the full Wilson pushforward.

\subsection{Fixed data and statement}
\label{fqc-setup}

Fix one allowed finite macrostep, including its preceding classical
depth $M$, current and output tori, ratios $R,J$, common angle radius
$b$, carrier exponent $a_0$, actual retained chart, and all strict
common-domain budgets. Use the finite activity quotients by equality
of vacuum-normalized totals on the full product ball. Their norms
are infima of anchored family norms.

Let $I$ be a compact interval contained in $(0,g_0]$, where $g_0$
meets the common-domain and two-input comparison bounds. For example,
take the continuous threshold
\begin{equation}
 T(g)=\max\{T_0,A\sqrt{\log(1/g)}\}.
 \label{fqc-threshold}
\end{equation}
More generally, any positive continuous threshold satisfying those
bounds suffices. Only continuity and boundedness on $I$ will be used.
Denote the actual finite quotient map by
\begin{equation}
 F_g(x)=T_{\mathrm{cl}}x+\Delta+\mathcal Q_{\det}
                         +\mathcal L_g(Q_x).
 \label{fqc-map}
\end{equation}
Here $T_{\mathrm{cl}}$ is classical background transport, the next two
terms are the fixed pivot and Gaussian determinant corrections, and
$Q_x$ is the inherited fluctuation difference. The stationary
backgrounds, precision and chart are independent of $g$. The
coupling is inserted \emph{inside} the nonlinear chart as
$\Phi_V(g\xi)$.

\begin{theorem}[Finite real-coupling continuity]
\label{fqc-theorem}
On the common admissible quotient ball, for instance
$\lVert x\rVert\leq\overline\epsilon/4$, the map
$(g,x)\mapsto F_g(x)$ is jointly continuous for $0<g\leq g_0$.
Both activity topologies use the same radius $b$ and exponent $a_0$.
Consequently the actual extracted update
\begin{equation}
 (\gamma,x)\longmapsto
 \left(\gamma-\beta_{\mathrm{out}}
                 (F_{\gamma^{-1/2}}(x)),
       P_{\mathrm{out}}F_{\gamma^{-1/2}}(x)\right)
 \label{fqc-extracted-map}
\end{equation}
is continuous on its admissible positive weak-coupling domain.
The finite maps $\beta_{\mathrm{out}},P_{\mathrm{out}}$ are normalized
by the actual finite reference and do not depend on $g$.
\end{theorem}

The proof permits constants depending exponentially on the finite
volume. It asserts neither a volume-uniform modulus nor continuity
at $g=0$. In particular it does not require a continuous choice of
near-minimal activity representatives.

\subsection{Strict inherited-field margin}
\label{fqc-margin-section}

Fix one input class $x$ and choose one admissible representative
$E$, held fixed as $g$ varies; one may take
$\lVert E\rVert<\overline\epsilon/2$.
The proved actual joint-family construction gives
\begin{equation}
 f_g(w,\xi)=\Omega_R(a_p(w),g\xi).
 \label{fqc-actual-field}
\end{equation}
Its absolute rooted family bound supplies, under the strict budgets,
\begin{equation}
 \lVert a_p(w)\rVert_\infty\leq b/4,\qquad
 \lVert f_g(w,\xi)\rVert_\infty\leq b/4
 \label{fqc-margin}
\end{equation}
for all $\lVert w\rVert_\infty<b$, $g\in I$, and every actual
retained link/color window of threshold $T(g)$. Fixed norm
conversions have already been included in the common-domain
constants. The argument below only requires a fixed inner radius
$sb$ with $s<1$.

Write $E_{\mathrm{raw}}(A)=\sum_H E_H(A|_H)$, and let
$V_{\mathrm{in}}$ be the number of input $J$-cells. The anchored norm
implies
\begin{equation}
 \sum_H\lVert E_H\rVert_b
       \leq V_{\mathrm{in}}\lVert E\rVert_{a_0;b}.
 \label{fqc-extensive-bound}
\end{equation}
Directional one-variable Cauchy on the product ball gives
\begin{equation}
 |E_{\mathrm{raw}}(A)-E_{\mathrm{raw}}(B)|
 \leq
 \frac{V_{\mathrm{in}}\lVert E\rVert_{a_0;b}}
      {(1-s)b}\,\lVert A-B\rVert_\infty,
 \qquad \lVert A\rVert_\infty,\lVert B\rVert_\infty\leq sb.
 \label{fqc-input-cauchy}
\end{equation}
Indeed the segment between $A,B$ stays in the inner product ball;
each unit max-norm direction has a complex disk of radius at least
$(1-s)b$. Apply Cauchy on any smaller disk, integrate the derivative
along the segment, sum using \eqref{fqc-extensive-bound}, and pass
to the limiting disk radius. Thus an arbitrary input need only be
bounded and holomorphic on its original $b$ ball. No extension
past that ball is assumed.

\subsection{Fixing the real integration window}
\label{fqc-integral-section}

Let $D$ be the total finite number of real retained scalar
coordinates. Let $\Xi_1\subset\mathbb R^D$ be the product of the
actual centered unit retained link/color balls. Its dilation
$T(g)\Xi_1$ is exactly the actual window; its boundary is
Lebesgue-null. Let $A(w)$ be the retained precision, with the
fixed analytic square-root determinant branch. Changing real
variables $\xi=T(g)\zeta$ gives the exact Gaussian expectation
\begin{align}
 Z_g(w)
 ={}&\frac{\det A(w)^{1/2}}{(2\pi)^{D/2}}\,T(g)^D
       \int_{\Xi_1}
       \exp\bigl\{
          -\tfrac12 T(g)^2\zeta^T A(w)\zeta
          +V_{\mathrm{base}}(w,T(g)\zeta;g)
 \notag\\[-2pt]
 &\hspace{39mm}
          +E_{\mathrm{raw}}(f_g(w,T(g)\zeta))
          -E_{\mathrm{raw}}(a_p(w))
       \bigr\}\,d\zeta .
 \label{fqc-integral}
\end{align}
The background-only pivot and determinant terms may remain outside
this expression as in \eqref{fqc-map}. Coupling powers, compact
gauge-insertion constants and flat Haar constants are separately
tracked vacuum factors. No contour deformation or replacement
by a whole-fiber Euclidean cutoff occurs. The chart argument is
$\Phi_V(gT(g)\zeta)$.

The independent classical domains used in constructing the actual
chart and base interaction give the following bounds on this
fixed regulator. The precision and its determinant branch are
bounded for all $\lVert w\rVert_\infty<b$. The local analytic
coefficient maps are bounded and uniformly Lipschitz in their
fluctuation arguments on a slightly larger fixed chart domain
containing every $gT(g)\zeta$. Moreover
$V_{\mathrm{base}}(w,T(g)\zeta;g)$ is uniformly bounded and
uniformly continuous in $g\in I$, uniformly over the full open
$b$ ball and $\zeta\in\Xi_1$.

The last assertion follows directly from the exact factors as
well. The classical and gauge Taylor remainders are analytic
functions of $g\xi$, divided by $g^2$ after removing the constant,
linear and quadratic terms. The remaining Haar and Jacobian
factors are analytic in the same argument. On a compact positive
$g$ interval all these expressions have bounded parameter
derivatives on the fixed strict chart margin. Their removable
singularities at zero coupling are not needed. An anchored bound
on a classical support series bounds its extensive sum at this
fixed volume.

Since $a_p$ is independent of $g$, continuity of $gT(g)$ and these
analytic bounds imply
\begin{equation}
 \sup_{\substack{\lVert w\rVert_\infty<b\\\zeta\in\Xi_1}}
 \bigl\lVert f_g(w,T(g)\zeta)
       -f_{g_*}(w,T(g_*)\zeta)\bigr\rVert_\infty
       \longrightarrow0.
 \label{fqc-field-continuity}
\end{equation}
Use \eqref{fqc-input-cauchy} for the inherited exponent. The
Gaussian term has the same uniform continuity because $A(w)$
is bounded and $T$ is continuous. Every exponent in
\eqref{fqc-integral} has a common finite absolute bound; the
prefactor is uniformly bounded and $\Xi_1$ has finite measure.
Uniform domination therefore proves
\begin{equation}
 \lVert Z_g-Z_{g_*}\rVert_{H^\infty(b)}
       \longrightarrow0 .
 \label{fqc-z-continuity}
\end{equation}
This is a supremum over the entire open $b$ ball, not merely
convergence on compact subballs.

\subsection{A lower bound on the full ball and logarithms}
\label{fqc-log-section}

The strict Gaussian theorem supplies nonvanishing throughout
the $b$ ball and a connected logarithm family with a common
anchored bound $H_I$ for $g\in I$. If $V_{\mathrm{out}}$ is the
number of output $J$-cells, its normalized total satisfies
\begin{equation}
 \sup_{\lVert w\rVert_\infty<b}
  |\log Z_g(w)-\log Z_g(0)|
       \leq 2V_{\mathrm{out}}H_I.
 \label{fqc-log-bound}
\end{equation}
The bound is unaffected by raw logarithm constants.
By \eqref{fqc-z-continuity} at $w=0$ and nonvanishing, on a
compact positive interval around $g_*$,
$m_0=\inf_{g\in I}|Z_g(0)|>0$. Hence
\begin{equation}
 \inf_{\substack{g\in I\\\lVert w\rVert_\infty<b}}
       |Z_g(w)|
 \geq m_0\exp(-2V_{\mathrm{out}}H_I)=:m_I>0.
 \label{fqc-lower-bound}
\end{equation}
Nonvanishing alone would not yield such an infimum on the open
ball; the uniform logarithm estimate is essential.

When $\lVert Z_g-Z_{g_*}\rVert_b<m_I/2$, use the convergent
power-series logarithm
\begin{equation}
 \ell_g(w)=
 \log\left(1+\frac{Z_g(w)-Z_{g_*}(w)}{Z_{g_*}(w)}\right),
 \qquad
 \lVert\ell_g\rVert_b\leq
          2m_I^{-1}\lVert Z_g-Z_{g_*}\rVert_b.
 \label{fqc-log-ratio}
\end{equation}
If $l_g$ is any selected holomorphic logarithm of $Z_g$, then
$l_g-l_{g_*}-\ell_g$ is constant on the connected product ball,
because its exponential is one. Thus
\begin{equation}
 \bigl\lVert[l_g-l_g(0)]-[l_{g_*}-l_{g_*}(0)]\bigr\rVert_b
 \leq 4m_I^{-1}\lVert Z_g-Z_{g_*}\rVert_b
 \longrightarrow0.
 \label{fqc-normalized-log-continuity}
\end{equation}
No continuity of a raw branch convention has been assumed.

\subsection{From totals to the finite quotient norm}
\label{fqc-quotient-section}

Let $H_{\mathrm{all}}$ be the full connected output $J$-cell
torus. Assign any bounded holomorphic normalized total $h$ to
this one carrier and zero to all others. This is an admissible
finite activity representation and gives
\begin{equation}
 \lVert[h]\rVert_{\mathrm{quot}}
       \leq e^{a_0V_{\mathrm{out}}}\lVert h\rVert_b.
 \label{fqc-single-carrier}
\end{equation}
If $h$ has invariant total, this representative retains that
property; coefficientwise gauge invariance is unnecessary.
The exterior terms in \eqref{fqc-map} are independent of $g$.
Applying \eqref{fqc-single-carrier} to
\eqref{fqc-normalized-log-continuity} proves, for fixed $x$,
\begin{equation}
 \lVert F_g(x)-F_{g_*}(x)\rVert_{\mathrm{quot}}
       \longrightarrow0.
 \label{fqc-fixed-input-continuity}
\end{equation}
The volume cost in \eqref{fqc-single-carrier} is intentional.
It avoids an unsupported inference about the norm continuity
of every canonical cluster coefficient or relocalized array.

\subsection{Joint continuity and finite histories}
\label{fqc-joint-section}

The finite two-input estimate is uniform on $I$:
\begin{equation}
 \lVert F_g(x)-F_g(y)\rVert_{\mathrm{quot}}
 \leq
 \left[B_T+C_{\mathrm{quant}}
                  \sup_{h\in I}\frac{hT(h)}b\right]
       \lVert x-y\rVert_{\mathrm{quot}}.
 \label{fqc-input-lipschitz}
\end{equation}
It follows from admissible near-minimal representatives
$Y,D,X=Y+D$ and the stronger comparison budget
$a_Q\overline\epsilon\leq\eta_*/(128e)$. It does not require
a continuous representative selection. For $(g,x)\to(g_*,x_*)$,
split the difference as
\[
 F_g(x)-F_{g_*}(x_*)
   =F_g(x)-F_g(x_*)+F_g(x_*)-F_{g_*}(x_*).
\]
Equations \eqref{fqc-input-lipschitz} and
\eqref{fqc-fixed-input-continuity} make both terms tend to
zero. Boundedness of the actual finite $\beta_{\mathrm{out}}$
and $P_{\mathrm{out}}$, together with continuity of
$\gamma\mapsto\gamma^{-1/2}$, proves
Theorem~\ref{fqc-theorem}.

Every fixed finite composition of the extracted updates is
continuous while all intermediate inputs and couplings remain
in their common admissible domains. This follows by induction
over the finite stages. One can work near a trajectory lying
strictly within these domains, or use the finite-induction
envelope to ensure admissibility on a prescribed parameter
interval.

Near-minimal representatives implementing the history may
jump with the initial coupling. Exact finite-total descent
makes the quotient update and its beta increment independent
of those choices. Raw vacuum factors remain separately
tracked by their exact compensation identities; an endpoint
statement about $\gamma$ does not require a continuous raw
representative convention.

Continuity alone supplies no endpoint crossing. An
intermediate-value argument also needs values on opposite
sides of the desired endpoint and control keeping the
intervening histories weak and admissible. A bound on the
magnitude of each increment may keep a sufficiently high
finite initial $\gamma$ weak, but it does not tune a fixed
terminal coupling or establish an ultraviolet trajectory.
The present conclusion remains finite-volume, positive-coupling,
and confined to the specified local analytic map.


% No preamble. Prefix dmt-. Exact actual pivot/Gaussian drift and trace kernels.
\section{An explicit trace formula for the pivot and Gaussian drift}
\label{dmt-section}

Fix the primitive ratio, a composite ratio \(R\), the actual
composite ordinary average, its selected pivots and the direct-\(R\)
internal contour gauge. The preceding homogeneous classical
depth \(M\) is arbitrary. All finite-regulator formulas in this
section are exact. Their infinite-volume interpretation uses one
chosen joint local-limit system with the previously proved uniform
exponential bounds. No numerical value, sign, old-depth limit or
universal beta coefficient is asserted.

\subsection{The exact external factors}
\label{dmt-definitions-section}

Use real Lie coordinates orthonormal for the normalized trace,
and their bilinear complexification. Write current link angles
as \(a\) and final coarse link angles as \(w\). Retained-coordinate
matrix traces use the actual Euclidean counting coordinates of
the Gaussian integral. They carry no extra old fine-volume factor.

Let \(F_R\) be the actual composite ordinary group average. Its
relative constraint and differential are
\begin{equation}
 F_a(B)=\frac1i\log\!\left[
     F_R(e^{iB}e^{ia})F_R(e^{ia})^{-1}\right],
 \qquad C_a=D_BF_a(0).
 \label{dmt-relative-constraint}
\end{equation}
For the unique selected pivot \(p(c)\) of a final bond \(c\),
\begin{equation}
 (C_a)_{c,p(c)}=R^{1-d}\pi_c(a),\qquad
 \pi_c(0)=I.
 \label{dmt-pivot-block}
\end{equation}
The actual pivot insertion, retained insertion and chart satisfy
\begin{align}
 C_ah_a&=I,\qquad
 K_a=E_{\rm free}-h_aC_aE_{\rm free},\notag\\
 D_z\Phi_a(0)&=K_a,\qquad
 D_z^2\Phi_a(0)[u,v]
       =-h_aD_B^2F_a(0)[K_au,K_av].
 \label{dmt-chart-identities}
\end{align}
The chart solves \(F_a(\Phi_a(z))=0\) and has retained
coordinates exactly \(z\).

Let \(G_R\) be the actual absolute contour gauge function.
It is not replaced by a background-relative penalty. Define
\begin{equation}
 \mathcal W_M(a)=S_M(e^{ia})+G_R(e^{ia}),\qquad
 f(a,z)=\frac1i\log[e^{i\Phi_a(z)}e^{ia}],
 \label{dmt-pulled-action}
\end{equation}
so \(f(a,0)=a\), and set
\begin{equation}
 \mathscr A_M(a)
    =D_z^2[\mathcal W_M(f(a,z))]_{z=0}.
 \label{dmt-independent-precision}
\end{equation}
This is the full analytic precision on the independent
background ball. Its Gaussian need not be centered at an
arbitrary independent background.

The actual physical branch \(a_*(w)\) satisfies
\[
 a_*(0)=0,\qquad F_R(e^{ia_*(w)})=e^{iw},\qquad
 P_R(e^{ia_*(w)};y,x)=I.
\]
It is the centered branch of the integral. Put
\begin{equation}
 A(w)=\mathscr A_M(a_*(w)),\qquad
 A_0=\mathscr A_M(0),\qquad C_0=A_0^{-1}.
 \label{dmt-physical-precision}
\end{equation}
The exact external logarithms are
\begin{equation}
 \Delta(w)=-\sum_c\log\det\pi_c(a_*(w)),\qquad
 \mathcal Q_{\det}(w)
       =-\frac12\log\det[A(w)A_0^{-1}],\qquad
 D_M=\Delta+\mathcal Q_{\det}.
 \label{dmt-leading-correction}
\end{equation}
All logarithmic branches in this display vanish at zero.
The factor \(R^{1-d}\) cancels between each pivot determinant
and its zero-background value.

The signs follow from the inverse pivot Jacobian and the
Gaussian factor \(\det A^{-1/2}\). There is no additional
background Haar factor \(j(a_*)\): the exact link Haar coordinates
are left-relative and give \(j(\Phi_a(g\xi))\), whose logarithm
vanishes at \(g=0\). The fluctuation straightening ratio is also
one at \(\xi=0\). Powers of \(g\), the gauge-insertion constant
and flat Gaussian constants are field independent and remain
outside \(D_M\). These definitions retain the actual mean,
composite map, chart and gauge; a direct-\(R\) replacement of
the constraint or an axial covariance would give a different
formula.

\subsection{The complete finite-regulator chain rule}
\label{dmt-chain-section}

For final coarse directions \(u,v\), write
\begin{equation}
 H_u=Da_*(0)u,\qquad B_{uv}=D^2a_*(0)[u,v].
 \label{dmt-stationary-responses}
\end{equation}
Define independent-background vertices by
\begin{align}
 \mathscr A_1[p]&=D\mathscr A_M(0)[p],&
 \mathscr A_2[p,q]&=D^2\mathscr A_M(0)[p,q],\notag\\
 \pi_{c,1}[p]&=D\pi_c(0)[p],&
 \pi_{c,2}[p,q]&=D^2\pi_c(0)[p,q].
 \label{dmt-background-vertices}
\end{align}
The physical derivatives are
\begin{align}
 A_u&=\mathscr A_1[H_u],&
 A_{uv}&=\mathscr A_2[H_u,H_v]+\mathscr A_1[B_{uv}],\notag\\
 \pi_{c,u}&=\pi_{c,1}[H_u],&
 \pi_{c,uv}&=\pi_{c,2}[H_u,H_v]+\pi_{c,1}[B_{uv}].
 \label{dmt-full-composition}
\end{align}

\begin{proposition}[Exact finite trace Hessian]
\label{dmt-finite-hessian}
For the actual finite integral,
\begin{align}
 D^2D_M(0)[u,v]
 ={}&-\sum_c\operatorname{tr}
       \{\pi_{c,uv}-\pi_{c,u}\pi_{c,v}\}\notag\\
 &-\frac12\operatorname{Tr}
       \{C_0A_{uv}-C_0A_uC_0A_v\}.
 \label{dmt-full-hessian}
\end{align}
Here lowercase trace acts on one pivot color block and
uppercase trace acts on all retained coordinates.
\end{proposition}
\begin{proof}
Differentiate the logarithmic determinants in
\eqref{dmt-leading-correction}, using
\[
 D\log\det A[B]=\operatorname{Tr}(A^{-1}B),
 \qquad D(A^{-1})[B]=-A^{-1}BA^{-1},
\]
and \(\pi_c(0)=I\). Products need not commute. Cyclicity of
the trace makes the answer symmetric in \(u,v\).
\end{proof}

In particular, constrained stationarity does not say that
the left current of \(S_M\) vanishes. It does not permit
deleting the constraint-curvature terms inside
\(\mathscr A_1\) or \(\mathscr A_2\).

\subsection{The second stationary response cancels after taking the scalar trace}
\label{dmt-color-section}

Introduce the independent-background scalar
\begin{equation}
 d_M(a)=-\sum_c\log\det\pi_c(a)
       -\frac12\log\det[\mathscr A_M(a)A_0^{-1}].
 \label{dmt-independent-scalar}
\end{equation}
This function is not asserted to be invariant under arbitrary
local background gauges. The constant-color symmetry suffices
for the next lemma.

\begin{lemma}[Semisimple cancellation of the first derivative]
\label{dmt-linear-cancellation}
For a semisimple Lie algebra,
\begin{equation}
                         Dd_M(0)=0.
 \label{dmt-zero-linear}
\end{equation}
Consequently the full Hessian in
\eqref{dmt-full-hessian} reduces to
\begin{align}
 D^2D_M(0)[u,v]
 ={}&-\sum_c\operatorname{tr}
 \{\pi_{c,2}[H_u,H_v]
       -\pi_{c,1}[H_u]\pi_{c,1}[H_v]\}\notag\\
 &-\frac12\operatorname{Tr}
 \{C_0\mathscr A_2[H_u,H_v]
       -C_0\mathscr A_1[H_u]C_0\mathscr A_1[H_v]\}.
 \label{dmt-reduced-hessian}
\end{align}
\end{lemma}
\begin{proof}
Let \(O=\operatorname{Ad}_k\) in the fixed real orthonormal
Lie basis. Endpoint covariance of the actual averages gives
\begin{equation}
                  \pi_c(Oa)=O\pi_c(a)O^{-1}.
 \label{dmt-pivot-color}
\end{equation}
The unique small pivot chart is equivariant because both
its constraint and retained-coordinate prescription are.
The absolute contour gauge is a sum of class functions,
and the actual scalar action is invariant under global
conjugation. Thus, with \(O_{\rm ret}\) the direct sum
on retained coordinates,
\begin{equation}
 \mathscr A_M(Oa)
       =O_{\rm ret}\mathscr A_M(a)O_{\rm ret}^{\,T},
 \qquad O_{\rm ret}^{\,T}O_{\rm ret}=I.
 \label{dmt-precision-color}
\end{equation}
The flat precision commutes with this representation.
Equations \eqref{dmt-pivot-color}--\eqref{dmt-precision-color}
prove invariance of \eqref{dmt-independent-scalar}, first
on the real domain and then by holomorphic continuation
with its zero-at-origin branches.

A semisimple Lie algebra has no invariant linear functional
on one adjoint color slot. Apply the invariance to a variation
on one current bond. Its linear functional is therefore zero.
Doing this for every bond proves \eqref{dmt-zero-linear}.
The argument applies separately to the pivot sum and the
Gaussian trace logarithm.

After substituting \eqref{dmt-full-composition} into
\eqref{dmt-full-hessian}, the coefficient of \(B_{uv}\)
is precisely \(Dd_M(0)[B_{uv}]\). This proves
\eqref{dmt-reduced-hessian}.
\end{proof}

This cancellation does not say that
\(\mathscr A_1=0\), \(\pi_{c,1}=0\), or that the current and
chart-curvature operators vanish. For an abelian factor,
retain \eqref{dmt-full-hessian} unless its linear functional
has been separately shown to vanish.

\subsection{The actual flat stationary lift}
\label{dmt-lift-section}

The first response appearing in
\eqref{dmt-reduced-hessian} is explicit in the proved flat
operators. Let
\begin{equation}
 C_R=D_a\{\log F_R(e^{ia})/i\}_{a=0},\qquad
 h=h_0,\quad K=K_0,\quad P_M=q_M+T_R^TT_R.
 \label{dmt-flat-data}
\end{equation}
Here \(q_M\) is the ordinary homogeneous classical Hessian,
and \(T_R\) is the linearized direct-\(R\) internal contour
map. The constraint in \(C_R\) is still the composite map.
The angle/counting normalization gives
\begin{equation}
 A_0=K^TP_MK,\qquad
 B_M=K^TP_Mh=K^Tq_Mh,\qquad
 H=h-KC_0B_M.
 \label{dmt-flat-lift}
\end{equation}
The simplification uses \(T_Rh=0\), since internal contours
use no pivots.

To verify the formula, use the identity-background affine
chart
\begin{align}
 \log F_R(e^{ia})/i&=C_Ra+\mathcal C_R(a),
 \qquad \mathcal C_R(0)=D\mathcal C_R(0)=0,\notag\\
 \mathfrak a(w,z)&=Kz+hw-h\mathcal C_R(\mathfrak a(w,z)).
 \label{dmt-affine-chart}
\end{align}
Put
\(\mathcal T_M(w,z)=\mathcal W_M(\mathfrak a(w,z))\).
Since the full action is stationary at zero,
\begin{equation}
 D_z\mathcal T_M(w,z)
       =A_0z+B_Mw+O((|w|+|z|)^2).
 \label{dmt-affine-stationarity}
\end{equation}
The positive flat retained Hessian and the established
local uniqueness identify its analytic critical point
with the physical branch. Thus
\(z_u=-C_0B_Mu\) and \(H_u=hu+Kz_u\), proving
\eqref{dmt-flat-lift}. No \(M^d\) or \(R^d\) transpose
factor belongs in these angle-coordinate matrices.

For completeness, the second response retained in the
general formula is also explicit. Set \(t_u=(u,z_u)\)
and define \(z_{uv}\) by the counting-pairing identity
\begin{equation}
 \langle\zeta,A_0z_{uv}\rangle
       =-D^3\mathcal T_M(0)[(0,\zeta),t_u,t_v]
       \quad\hbox{for every retained }\zeta.
 \label{dmt-second-stationary}
\end{equation}
Then
\begin{equation}
 B_{uv}=-hD^2\mathcal C_R(0)[H_u,H_v]+Kz_{uv}.
 \label{dmt-second-background}
\end{equation}
These equations follow by twice differentiating the retained
stationarity equation and \eqref{dmt-affine-chart}; they are
available even when the color cancellation is not.

\subsection{The complete precision vertices}
\label{dmt-precision-vertices-section}

Set
\begin{equation}
 E_a=D_zf(a,0),\qquad T_a=D_z^2f(a,0).
 \label{dmt-chart-tensors}
\end{equation}
With
\[
 J_aX=\int_0^1\operatorname{Ad}(e^{ita})X\,dt,\qquad
 \Xi_a(B)=\frac1i\log[e^{iB}e^{ia}],
\]
the exact tensors are
\begin{align}
 E_a&=J_a^{-1}K_a,\notag\\
 T_a[u,v]&=-J_a^{-1}h_aD_B^2F_a(0)[K_au,K_av]
       +D_B^2\Xi_a(0)[K_au,K_av].
 \label{dmt-both-chart-curvatures}
\end{align}
Consequently the arbitrary-background precision is
\begin{equation}
 \mathscr A_M(a)[u,v]
       =D^2\mathcal W_M(a)[E_au,E_av]
                  +D\mathcal W_M(a)[T_a(u,v)].
 \label{dmt-complete-precision}
\end{equation}
The two terms of \(T_a\) are respectively nonlinear constraint
curvature and conversion from the left-relative exponential
coordinate to the ordinary link logarithm. The absolute gauge
function remains inside \(\mathcal W_M\), so its nonphysical
background curvature is included as well. Its gradient
vanishes on the physical branch, but this does not permit
replacing the independent analytic family by a relative
quadratic gauge form.

Write \(W_n=D^n\mathcal W_M(0)\), so that
\(W_1=0\), \(W_2=P_M\), \(E_0=K\). Define
\[
 E_1[p]=D_aE_0[p],\qquad E_2[p,q]=D_a^2E_0[p,q],
 \qquad T_0=T_{a=0},\qquad T_1[p]=D_aT_0[p].
\]
For retained directions \(r,s\), direct differentiation gives
\begin{align}
 \mathscr A_1[p][r,s]
 ={}&W_3[p,Kr,Ks]
       +W_2[E_1[p]r,Ks]+W_2[Kr,E_1[p]s]\notag\\
 &+W_2[p,T_0(r,s)].
 \label{dmt-first-precision-vertex}
\end{align}
The last term is a first-background-order current and
constraint contribution; it remains despite \(W_1=0\).
The second vertex is
\begin{align}
 \mathscr A_2[p,q][r,s]
 ={}&W_4[p,q,Kr,Ks]\notag\\
 &+W_3[p,E_1[q]r,Ks]+W_3[p,Kr,E_1[q]s]\notag\\
 &+W_3[q,E_1[p]r,Ks]+W_3[q,Kr,E_1[p]s]\notag\\
 &+W_2[E_2[p,q]r,Ks]+W_2[Kr,E_2[p,q]s]\notag\\
 &+W_2[E_1[p]r,E_1[q]s]
       +W_2[E_1[q]r,E_1[p]s]\notag\\
 &+W_3[p,q,T_0(r,s)]\notag\\
 &+W_2[p,T_1[q](r,s)]
       +W_2[q,T_1[p](r,s)].
 \label{dmt-second-precision-vertex}
\end{align}
Only the \(W_1D_a^2T\) term disappears at zero.
The pivot vertices follow by twice differentiating
\begin{equation}
             \pi_c(a)=R^{d-1}(D_BF_a(0))_{c,p(c)}.
 \label{dmt-pivot-vertex-definition}
\end{equation}
All chart operations here use the actual finite-dimensional
logarithm and mean, and have finite range at the fixed \(R\).

\subsection{The classical cubic and quartic vertices}
\label{dmt-classical-vertices-section}

The classical vertices in
\eqref{dmt-first-precision-vertex}--\eqref{dmt-second-precision-vertex}
can be evaluated from the previously proved simultaneous
corrected-slice equations. This avoids an unproved bound for
the ordinary rooted average's physical adjoint.

Use the fine extension \(\mathcal H_M\), constrained Green
operator \(\mathcal C_M\), flat coarse Hessian \(q_M\),
corrected constraint remainder \(\mathcal R_M\), Wilson
jet \(f_\eta\), and its finite jet map \(\mathcal S\).
With the physical fine transposes and measures, the exact
pair \(Y=(A,m)\) obeys
\begin{align}
 A={}&\mathcal H_Ma-\mathcal H_M\mathcal R_M(A)
       -\mathcal C_M\mathcal S^Tf_\eta(\mathcal SA)
       +\mathcal C_MD\mathcal R_M(A)^Tm,\notag\\
 m={}&q_Ma-q_M\mathcal R_M(A)
       +(\mathcal S\mathcal H_M)^Tf_\eta(\mathcal SA)
       -\mathcal H_M^TD\mathcal R_M(A)^Tm.
 \label{dmt-classical-pair}
\end{align}
Write \(Y=\mathcal L_Ma+\mathcal V_M(Y)\), where
\(\mathcal L_M=(\mathcal H_M,q_M)\) and
\(\mathcal V_M(0)=D\mathcal V_M(0)=0\).
For \(\mathcal V_n=D^n\mathcal V_M(0)\), its responses are
\begin{align}
 Y_1[u]&=\mathcal L_Mu,\notag\\
 Y_2[u,v]&=\mathcal V_2[\mathcal L_Mu,\mathcal L_Mv],\notag\\
 Y_3[u,v,t]&=\mathcal V_3[
                  \mathcal L_Mu,\mathcal L_Mv,\mathcal L_Mt]\notag\\
 &\quad+\mathcal V_2[Y_2[u,v],\mathcal L_Mt]
       +\mathcal V_2[Y_2[u,t],\mathcal L_Mv]
       +\mathcal V_2[Y_2[v,t],\mathcal L_Mu].
 \label{dmt-classical-response-recursion}
\end{align}
These are derivatives, not Taylor coefficients, so no factorial
is omitted. The scalar envelope \(DS_M(a)=m(a)\) and equality
of the corrected and ordinary action values yield
\begin{align}
 D^3S_M(0)[u,v,t]
       &=\langle(Y_2[u,v])_m,t\rangle,\notag\\
 D^4S_M(0)[u,v,t,r]
       &=\langle(Y_3[u,v,t])_m,r\rangle.
 \label{dmt-classical-scalar-vertices}
\end{align}
Add the explicit local derivatives of \(G_R\) to obtain
\(W_3,W_4\).

These are finite compositions of the proved propagators and
finite-range analytic vertices. The uniform estimates use
complete physical cell spaces and smooth supported
reconstruction, not unweighted sums over \(M^4\) microscopic
coordinates. The fine \(\eta^4\) measure is already in the
physical transposes of \eqref{dmt-classical-pair}. It must
not be inserted again into the retained-coordinate traces.

\subsection{Absolute trace convergence and local kernels}
\label{dmt-convergence-section}

The strict precision theorem supplies symmetric pieces with
matrix rows and columns in their carrier and an analytic bound
of the following form, at fixed positive
\(\sigma,a,\rho\):
\begin{equation}
 \sup_i\sum_{Y\ni i}\sum_{j\in Y}
 e^{\sigma\tau(Y)+a|Y|}
 \|(\mathscr A_M-A_0)_{Y,ij}\|_\rho\le N.
 \label{dmt-strict-row-input}
\end{equation}
This is uniform in preceding depth and permitted volume.
The precision-composition theorem supplies the same form
after stationary substitution, keeping the actual matrix
rows rooted in their original carrier. The flat covariance
has an exponential Schur bound. The pivot scalar and
stationary field have their established strict and rooted
bounds.

\begin{lemma}[Source-rooted trace bounds]
\label{dmt-trace-convergence}
On a fixed smaller analytic domain, the actual Hessian
kernel
\[
 \mathcal K_{xy}=D^2D_M(0)[e_x,e_y]
\]
has an exponentially summable row bound
\begin{equation}
       \sup_x\sum_y e^{m d(x,y)}|\mathcal K_{xy}|
                       \le C_D
 \label{dmt-kernel-decay}
\end{equation}
for some \(m>0\). The traces in both
\eqref{dmt-full-hessian} and
\eqref{dmt-reduced-hessian} are absolutely convergent on
compact or summable variation directions. Constants may
depend on the fixed geometry, color and radii, but not
on preceding depth or volume.
\end{lemma}
\begin{proof}
A derivative in one background coordinate \(x\) acts only
on carriers containing \(x\). Cauchy bounds it by a fixed
constant times \(\rho^{-1}\) times the original majorant.
Two derivatives require both source slots in the carrier
and give \(\rho^{-2}\).

Reduce the spatial exponent in
\eqref{dmt-strict-row-input}. For each matrix row \(i\),
the unused exponent pays \(e^{-\delta d(i,x)}\), since
\(\tau(Y)\ge d(i,x)\) when \(i,x\in Y\). Its sum over
\(i\) is the finite lattice bound \(S_\delta\), also
uniformly on periodic graphs. With a second source \(y\),
spend another exponent portion on
\(\tau(Y)\ge d(x,y)\).

For the composed precision, this yields, after fixed
exponent reductions,
\begin{align}
 \sum_{ij}e^{m_1[d(i,x)+d(j,x)]}
                   \|(A_x)_{ij}\|&\le C_1,\notag\\
 \sum_{ij}\|(A_{xy})_{ij}\|
                   &\le C_2e^{-m_1d(x,y)}.
 \label{dmt-entry-source-bounds}
\end{align}
For the first line, use
\(d(i,x)+d(j,x)\le2\tau(Y)\) and reserve an additional
exponent for the row summation. Finite fiber and color
factors are absorbed into the constants.

The same argument before composition, followed by the
exponentially decaying external lift \(H\), bounds the
separate terms of \eqref{dmt-reduced-hessian}. Thus the
proof does not use cancellation of non-absolutely
convergent quantities.

The trace with \(A_{xy}\) is bounded by its absolute-entry
sum and the covariance entry envelope. For the other
trace expand
\begin{equation}
 \operatorname{Tr}(C_0A_xC_0A_y)
 =\sum_{i,j,k,l}\operatorname{tr}
     \{C_{0,ij}(A_x)_{jk}C_{0,kl}(A_y)_{li}\}.
 \label{dmt-trace-expansion}
\end{equation}
The source-root bounds and covariance envelope make this
absolutely convergent. The \(x\)-root weight at \(j\),
the covariance weight between \(j\) and \(i\), and the
\(y\)-root weight at \(i\) jointly pay \(d(x,y)\).
All remaining factors are absolutely summable. Hence
\[
 |\operatorname{Tr}(C_0A_xC_0A_y)|
                         \le C_3e^{-m_2d(x,y)}
\]
for a smaller positive \(m_2\). The pivot traces have the
same bound by local dependence before composition and the
rooted stationary family, or directly by their composed
strict scalar expansion. A final exponent reduction
allows the row sum in \eqref{dmt-kernel-decay}.
\end{proof}

Grouping retained coordinates into auxiliary \(J\)-cells
only regroups these traces; it inserts no extra factor
\(J^d\) into the exact identities.

On the infinite lattice a total log determinant at a
noncompact background need not exist as an extensive
number. The displayed formulas concern local kernels
and derivatives on compact or summable directions.
They are defined by the absolutely convergent sums above
or by a chosen joint local limit of the actual finite
matrices and vertices. The common majorants permit that
limit. No arbitrary finite activity is identified with
the periodization of an independently selected infinite
one.

\subsection{A convergent Maxwell moment of the traces}
\label{dmt-Maxwell-section}

Assume the chosen infinite total has the established
translation, hypercubic and local gauge identities, and
fix one simple color factor. The midpoint Hessian symbol
then has
\begin{equation}
 \widehat{\mathcal K}_{\mu\nu}^{ab}(p)
 =\delta^{ab}\beta(D_M)
       (|p|^2\delta_{\mu\nu}-p_\mu p_\nu)+O(|p|^4).
 \label{dmt-Maxwell-symbol}
\end{equation}
For \(\mu\ne\nu\), choose one trace-orthonormal color
vector \(t^a\), and let \(e_{x,\mu,a}\) be its final
unit-spacing link direction. The midpoint offsets
cancel for the diagonal orientation \(\mu,\mu\).
Therefore
\begin{equation}
 \boxed{\displaystyle
 \beta(D_M)=-\frac12\sum_{r\in\mathbb Z^4}r_\nu^2
       D^2D_M(0)[e_{0,\mu,a},e_{r,\mu,a}].}
 \label{dmt-Maxwell-moment}
\end{equation}
There is no color sum here. Hypercubic and color symmetry
make the choices equivalent. The minus sign comes from
twice differentiating \(e^{ip\cdot r}\); the factor
\(1/2\) extracts the coefficient of \(p_\nu^2\).
The actual classical reference has coefficient one in
this normalization.

Writing \(H_r=H e_{r,\mu,a}\), substitution of
\eqref{dmt-reduced-hessian} gives
\begin{align}
 \beta(D_M)
 ={}&\frac12\sum_r r_\nu^2\sum_c\operatorname{tr}
 \{\pi_{c,2}[H_0,H_r]
       -\pi_{c,1}[H_0]\pi_{c,1}[H_r]\}\notag\\
 &+\frac14\sum_r r_\nu^2\operatorname{Tr}
 \{C_0\mathscr A_2[H_0,H_r]
       -C_0\mathscr A_1[H_0]C_0\mathscr A_1[H_r]\}.
 \label{dmt-explicit-Maxwell-trace}
\end{align}
Every sum is absolutely convergent by
Lemma~\ref{dmt-trace-convergence}. Truncating the external
displacement at distance \(L\) has error at most
\(C_{m'}e^{-m'L}\) for any \(0<m'<m\).
Internal trace tails likewise have exponential bounds
after anchoring the sources and spending a spare exponent.

At an actual finite regulator, the trace Hessian formulas
remain exact. Apply the chosen symmetric short-lift
moment and divide by the moment of the actual finite
reference to obtain its finite Maxwell ratio. That
reference denominator is \(1+O(e^{-cN})\), not
automatically exactly one. No finite Ward identity at
discrete momenta is interpolated to obtain
\eqref{dmt-Maxwell-symbol}; its use requires the stated
infinite kernel and symmetries, or the separately
justified local-limit passage.

\subsection{What has and has not been reduced}
\label{dmt-scope-section}

The leading zero-input drift is expressed by actual local
pivot vertices, the flat stationary lift
\eqref{dmt-flat-lift}, the retained covariance, and the
full precision vertices
\eqref{dmt-first-precision-vertex}--\eqref{dmt-second-precision-vertex}.
The classical cubic and quartic vertices have the finite
operator recursion
\eqref{dmt-classical-pair}--\eqref{dmt-classical-scalar-vertices}.
All signs and coordinate normalizations are fixed.

The second physical stationary response was included
before being removed by the proved semisimple color
identity. The physical current was not replaced by an
independent parameter, and neither chart-curvature
contribution was omitted.

Evaluating \eqref{dmt-explicit-Maxwell-trace}, comparing
its values as \(M\) grows, determining a sign or a
universal coefficient, and controlling a physical
coupling trajectory remain further tasks. Absolute
convergence of the trace expression alone proves none
of those conclusions.


% No preamble. Prefix dgb-. Exact finite identities, not a uniform comparison.
\section{An exact gauge and coarea bridge for the actual determinant}
\label{dgb-section}

This section rewrites the actual leading correction as coupled
vector, ghost and coarse-constraint determinants. It keeps the
constraint-current curvature and every finite orbit-density factor.
It also exhibits a soft mode that prevents a separate uniformly
massive estimate for the based ghost. The bounded finite-part
comparison with a standard background-Feynman determinant remains
a separate estimate.

\subsection{The actual multiplier Hessian}
\label{dgb-setup-section}

Fix a finite compatible regulator and endpoint $K=\ell^k$.
Let $U=V_*^{M,K}(W)$ be the physical stationary representative,
so $F_K(U)=W$ and $P_K(U)=I$. All matrices use the actual
current-link angle and scalar-site counting coordinates in a real
trace-orthonormal color basis.

In left-relative coordinates define
\begin{equation}
 s(x)=S_M(e^{ix}U),\qquad
 f(x)=\frac1i\log\!\left[F_K(e^{ix}U)W^{-1}\right],
 \qquad C=f'(0).
 \label{dgb-relative-coordinates}
\end{equation}
The target coordinate is a relative increment at $W$, not its
ordinary logarithm. Constrained stationarity gives a unique
multiplier $\lambda$ with $s'(0)=C^T\lambda$. Put
\begin{equation}
 \mathcal L=s''(0)-\sum_c\lambda_c f_c''(0).
 \label{dgb-Lagrangian}
\end{equation}
This is the Hessian of the critical Lagrangian
$s-\lambda\cdot f$. Its second term is the full
constraint-current curvature.

Let $C_p$ be the invertible actual pivot block, $I_p$ the pivot
insertion, and $E$ the retained insertion. Write
\begin{equation}
 h=I_pC_p^{-1},\qquad K_0=E-hCE,\qquad
 Ch=I,\qquad CK_0=0.
 \label{dgb-chart-linear}
\end{equation}
The subscript distinguishes $K_0$ from the endpoint ratio $K$.
The nonlinear source chart satisfies
\begin{equation}
 \Phi'(0)=K_0,\qquad
 \Phi''(0)[u,v]=-h f''(0)[K_0u,K_0v].
 \label{dgb-chart-curvature}
\end{equation}
Consequently
\[
 D_z^2s(\Phi(z))\big|_{z=0}=K_0^T\mathcal L K_0 .
\]
The actual absolute contour penalty vanishes at the saddle;
its Hessian is $T^TT$, where $T$ differentiates the contour
logarithms in the established normalization. Thus
\begin{equation}
 A_T=K_0^T(\mathcal L+T^TT)K_0>0.
 \label{dgb-actual-precision}
\end{equation}
Apart from field-independent constants, the actual leading
amplitude is
\begin{equation}
                  |\det C_p|^{-1}\det(A_T)^{-1/2}.
 \label{dgb-actual-amplitude}
\end{equation}
The real determinant signs are fixed from the flat convention.
All logarithms below are normalized at $W=1$ and then continued
on their chosen holomorphic branches.

\subsection{The based gauge nullspace}
\label{dgb-nullspace-section}

Let $D_b$ send scalar site parameters vanishing at every
$K$-center to their left-relative infinitesimal gauge variations
at $U$. Exact based invariance of the ordinary constraint gives
\begin{equation}
                              CD_b=0.
 \label{dgb-based-constraint}
\end{equation}
The scalar $s-\lambda\cdot f$ is based-gauge invariant and is
critical at zero. Differentiating its gauge identity there gives
\begin{equation}
                              \mathcal L D_b=0.
 \label{dgb-based-null}
\end{equation}
This need not hold for $s''$ without its multiplier-curvature
term. The actual contour endpoint covariance gives
\begin{equation}
                              TD_b=-I
 \label{dgb-contour-gauge}
\end{equation}
with the standard orientation. Changing the generator sign
changes only this sign. A fixed color normalization in the
gauge penalty contributes only a background-independent
factor, canceled by normalization.

The based action is free, hence $D_b$ is injective and
\begin{equation}
                              G_b=D_b^TD_b>0.
 \label{dgb-based-ghost}
\end{equation}
On $\ker C$, the form $\mathcal L$ is nonnegative, has
nullspace precisely $\operatorname{ran}D_b$, and is positive
on its quotient. Indeed, for any $x\in\ker C$, use
\eqref{dgb-contour-gauge} to add a unique based gauge vector
making $Tx=0$. Equation \eqref{dgb-based-null} leaves its
$\mathcal L$ value unchanged. On this contour-zero complement,
the positive form \eqref{dgb-actual-precision} equals
$\mathcal L$. Equality occurs only for a vertical vector.

\subsection{A finite Gaussian gauge-change identity}
\label{dgb-gauge-change-section}

\begin{lemma}[Gauge replacement with its orbit density]
\label{dgb-gauge-lemma}
Let $B_0$ be a symmetric nonnegative form on a finite space
$Z$, with nullspace $\operatorname{ran}V$, positive on the
quotient. Suppose $F_1,F_2:Z\to\mathbb R^s$ have invertible
$F_jV$. Then
\begin{equation}
 |\det(F_1V)|\det(B_0+F_1^TF_1)^{-1/2}
   =|\det(F_2V)|\det(B_0+F_2^TF_2)^{-1/2}.
 \label{dgb-gauge-identity}
\end{equation}
\end{lemma}

\begin{proof}
Choose one complement $E_0$ to $\operatorname{ran}V$ and set
$z=E_0u+V\alpha$. The change-of-variable Jacobian is the same
in the two integrals. The $B_0$ term depends only on $u$.
Integration over $\alpha$ yields
$(2\pi)^{s/2}/|\det(F_jV)|$, independently of $u$, by
translation and an invertible linear change of variables.
The remaining integral in $u$ is common. This proves
\eqref{dgb-gauge-identity}; no orbit-coordinate density
has been discarded.
\end{proof}

Apply the lemma on retained coordinates, with vertical insertion
$V$ characterized by $K_0V=D_b$. Take
$F_1=TK_0$ and $F_2=D_b^TK_0$. Equations
\eqref{dgb-contour-gauge} and \eqref{dgb-based-ghost} give
$F_1V=-I$ and $F_2V=G_b$. Therefore
\begin{equation}
 \det(A_T)^{-1/2}
     =\det(G_b)\det(A_b)^{-1/2},\qquad
 A_b=K_0^T(\mathcal L+D_bD_b^T)K_0.
 \label{dgb-divergence-replacement}
\end{equation}
This replaces the actual contour gauge by based covariant
divergence. It does not yet replace it by a full-site massive
background-Feynman gauge.

\subsection{Coarea removes the pivot}
\label{dgb-coarea-section}

Choose any symmetric coarse matrix $B$ such that
\begin{equation}
 Q_b=\mathcal L+D_bD_b^T+C^TBC>0,\qquad
 M_b=CQ_b^{-1}C^T>0.
 \label{dgb-completion}
\end{equation}
A completion exists at each fixed finite saddle: the first
two terms are positive on $\ker C$, so a sufficiently large
scalar $B$ makes the complementary Schur block positive.
The condition is positivity of $Q_b$. Merely requiring
$B>0$ need not suffice; conversely $B$ itself may be indefinite.

Let $n$ and $m$ be the ambient and constraint dimensions.
The same constrained Gaussian has the exact evaluations
\begin{align}
 \int_{\mathbb R^n}\delta(Cx)e^{-x^TQ_bx/2}\,dx
  &=(2\pi)^{(n-m)/2}|\det C_p|^{-1}\det(A_b)^{-1/2}
       \notag\\
  &=(2\pi)^{(n-m)/2}
           [\det Q_b\,\det M_b]^{-1/2}.
 \label{dgb-coarea-integral}
\end{align}
For the first, write $x=K_0z+hy$. In retained/pivot order,
this transformation has determinant $(\det C_p)^{-1}$,
and $Cx=y$. For the second, the normalized full Gaussian
of covariance $Q_b^{-1}$ sends $Cx$ to a Gaussian of covariance
$M_b$; its density at zero gives the factor.
Equivalently,
\begin{equation}
 \det Q_b\,\det(CQ_b^{-1}C^T)
                         =(\det C_p)^2\det A_b .
 \label{dgb-coarea-determinant}
\end{equation}

The delta in \eqref{dgb-coarea-integral} uses the target's
relative Lie coordinate. The actual target measure is group
Haar, whose density at that coordinate origin is the fixed
$h_0$. It supplies no factor $j(\log W)$.
If one instead uses induced Euclidean measure on $\ker C$,
its coarea factor is precisely included in
\eqref{dgb-coarea-integral}.

For a scalar expression $F$, write
$\mathcal N_0F=F(W)-F(1)$. Combining
\eqref{dgb-actual-amplitude},
\eqref{dgb-divergence-replacement} and
\eqref{dgb-coarea-integral} proves the exact identity
\begin{equation}
 \boxed{\displaystyle
 D_{M,K}
   =\mathcal N_0\left[
       \log\det G_b-\frac12\log\det Q_b
                    -\frac12\log\det M_b\right].}
 \label{dgb-main-identity}
\end{equation}
The pivot has disappeared, but the coarse-constraint
determinant remains. Every completion allowed in
\eqref{dgb-completion} gives the same combined expression.
The individual terms need not be completion-independent
or locally gauge invariant.

At fixed regulator these are identities on a positive real
germ and then on their chosen analytic branches. The algebraic
identity \eqref{dgb-coarea-determinant} also continues wherever
its inverses exist. No uniform complex domain for the three
individual logarithms, as $K$ or volume grows, follows from
the combined identity.

\subsection{A positive completion from the actual minimum action}
\label{dgb-residual-section}

In the same relative target coordinate set
\[
 s_{\rm eff}(y)=S_{MK}(e^{iy}W),\qquad
 q_{\rm eff}=s_{\rm eff}''(0).
\]
The envelope identity gives $s_{\rm eff}'(0)=\lambda$.
The residual scalar
\[
                 \mathcal R(x)=s(x)-s_{\rm eff}(f(x))
\]
is critical and nonnegative near zero by the actual local
minimum comparison. It is invariant under the full local
fine gauge group, because both scalar actions are locally
gauge invariant and the composite map is equivariant.
Its Hessian is therefore
\begin{equation}
 H_{\rm res}=\mathcal L-C^Tq_{\rm eff}C\geq0,\qquad
 H_{\rm res}D_f=0
 \label{dgb-residual-Hessian}
\end{equation}
for every fine infinitesimal gauge direction, with any
convenient fixed global base convention.
The nonnegativity uses the proved minimum comparison,
not stationarity alone.

Let $H$ be the actual derivative of the minimum section in
these relative coordinates. Since $CH=I$ and the residual
vanishes identically along that section, its positive
semidefinite Hessian obeys
\begin{equation}
                     H_{\rm res}H=0,\qquad
                     \mathcal LH=C^Tq_{\rm eff}.
 \label{dgb-residual-lift}
\end{equation}
No independently chosen flat lift is substituted for $H$.

A canonical completion is
\begin{equation}
                Q_b=H_{\rm res}+D_bD_b^T+C^TC>0.
 \label{dgb-canonical-completion}
\end{equation}
Indeed a null vector would satisfy $Cx=0$, $D_b^Tx=0$
and $x^T\mathcal Lx=0$. The constrained quotient positivity
then gives $x=D_b\alpha$, and $G_b\alpha=0$ implies $x=0$.
This is \eqref{dgb-completion} with $B=I-q_{\rm eff}$,
which need not be positive itself. No uniform lower
spectral bound for \eqref{dgb-canonical-completion} is
asserted as $K$ grows.

The general coarse compression can also be identified
exactly. Set
\begin{equation}
                   H_b=H-D_bG_b^{-1}D_b^TH .
 \label{dgb-adjusted-lift}
\end{equation}
Then $CH_b=I$, $D_b^TH_b=0$, $H_{\rm res}H_b=0$, and
$\mathcal LH_b=C^Tq_{\rm eff}$. For every completion
\eqref{dgb-completion}, it follows that
\begin{equation}
 Q_bH_b=C^T(q_{\rm eff}+B),\qquad
 q_{\rm eff}+B=H_b^TQ_bH_b>0.
 \label{dgb-effective-Schur}
\end{equation}
Multiplication by the inverses gives
\begin{equation}
 Q_b^{-1}C^T=H_b(q_{\rm eff}+B)^{-1},\qquad
 \boxed{\,M_b=(q_{\rm eff}+B)^{-1}\,}.
 \label{dgb-coarse-identification}
\end{equation}
Thus \eqref{dgb-main-identity} is equivalently
\begin{equation}
 D_{M,K}=\mathcal N_0\left[
      \log\det G_b-\frac12\log\det Q_b
                   +\frac12\log\det(q_{\rm eff}+B)\right].
 \label{dgb-effective-determinant}
\end{equation}
The coarse covariance is the inverse of the actual
relative-coordinate effective Hessian with its completion.
For the canonical choice $B=I-q_{\rm eff}$, it is exactly
$M_b=I$ and that coarse logarithm vanishes. The coarse
physics remains inside
$H_{\rm res}=\mathcal L-C^Tq_{\rm eff}C$ and hence in $Q_b$;
no uniform reference comparison is thereby obtained.

\subsection{A soft based-ghost mode}
\label{dgb-soft-ghost-section}

At the flat background $D_b$ is the nearest-neighbor gauge
gradient on scalar functions vanishing at the $K$-centers.
Take $K\geq3$ on a compatible unit-current torus and choose
one real color. Let $\varphi$ be zero at all centers and
one elsewhere. If there are $n_c$ centers, there are
$n_f=K^dn_c$ current sites, and exactly $2d$ incident
unit-difference bonds per center. Hence
\begin{equation}
 \|\varphi\|_2^2=(K^d-1)n_c,\qquad
 \|D_b\varphi\|_2^2=2dn_c,\qquad
 \lambda_{\min}(G_b)\leq\frac{2d}{K^d-1}.
 \label{dgb-Rayleigh}
\end{equation}
This is an upper Rayleigh bound, not a matching spectral
asymptotic. In four dimensions it implies
\[
                         \|G_b^{-1}\|\geq\frac{K^4-1}{8}.
\]
A final-physical-scale mass is $K^{-2}$ in the current
coordinates. Thus the estimate $\|G_b^{-1}\|\leq CK^2$
is false uniformly in $K$. After resetting final spacing
to one, $K^2G_b$ still has a Rayleigh eigenvalue tending
to zero.

This counterexample does not refute a bounded Maxwell
finite-part difference for the complete expression
\eqref{dgb-main-identity}. It rules out estimating the
based ghost separately as a uniformly massive operator.
The vector and coarse determinant must be kept coupled.

\subsection{An explicit full-site and coarse Schur decomposition}
\label{dgb-full-Schur-section}

Fix one final center as a global base. Let $D_f$ now denote
the fine gauge generator for all scalar sites except that
base, and split its domain into the non-$K$-centers and the
other $K$-centers:
\[
                            D_f=[\,D_b,D_e\,].
\]
These are disjoint site-coordinate columns, not orthogonal
images. The globally based action is free. Therefore
\begin{align}
 G_f&=D_f^TD_f>0,\notag\\
 S_g&=D_e^T(I-D_bG_b^{-1}D_b^T)D_e>0,\notag\\
 \det G_f&=\det G_b\,\det S_g .
 \label{dgb-ghost-Schur}
\end{align}
For any positive $Q_b$ as above, set
\begin{equation}
 Q_f=Q_b+D_eD_e^T,\qquad
 S_v=I+D_e^TQ_b^{-1}D_e.
 \label{dgb-vector-Schur-data}
\end{equation}
The determinant lemma gives
\begin{equation}
                         \det Q_f=\det Q_b\,\det S_v .
 \label{dgb-vector-Schur}
\end{equation}
Consequently the actual correction is exactly
\begin{equation}
 D_{M,K}=\mathcal N_0\left[
   \log\det G_f-\frac12\log\det Q_f
   -\log\det S_g+\frac12\log\det S_v
   -\frac12\log\det M_b\right].
 \label{dgb-full-decomposition}
\end{equation}
$S_g,S_v$ are coarse-site matrices; $M_b$ is a coarse-bond
matrix. The first two terms resemble a vector/ghost
combination, but all coarse corrections remain.
$G_f$ is globally based; it is not yet the translation-invariant
all-site massive ghost of a reference calculation.
If no nonbase coarse site remains, the corresponding
zero-dimensional determinants are one. For the canonical
completion $M_b=I$ by \eqref{dgb-coarse-identification};
the general formula retains all three coarse terms.

For \eqref{dgb-canonical-completion}, the vector operator is
\begin{equation}
 Q_f=s''+D_fD_f^T-\lambda\cdot f''
                         +C^T(I-q_{\rm eff})C.
 \label{dgb-vector-operator}
\end{equation}
If a proposed matched massive reference is
$P_\eta=s''+D_fD_f^T+\eta^2I$, with $\eta=K^{-1}$,
the exact difference is
\begin{equation}
 Q_f-P_\eta=-\lambda\cdot f''
                     +C^T(I-q_{\rm eff})C-\eta^2I.
 \label{dgb-reference-difference}
\end{equation}
The multiplier term vanishes at the flat background,
but its first and second background derivatives can
contribute to the determinant Hessian. Neither that
vanishing nor the coarse rank bounds its Maxwell
coefficient by itself. Equation
\eqref{dgb-reference-difference} is a comparison interface,
not an identification with an operator named in a
different background-field convention.
Any logarithmic comparison using $P_\eta$ additionally
requires its invertibility on the chosen finite germ.
No uniform such neighborhood is inferred here.

\subsection{A uniformly bounded coarse effective-Hessian factor}
\label{dgb-coarse-bound-section}

The general identification in
\eqref{dgb-coarse-identification} permits a uniform estimate
for one actual comparison term. Use the established strict
scalar family of the entire action $S_{MK}$, including its
pair-supported quadratic baseline. On a fixed small complex
ball it has uniform spatial bounds and an arbitrary
prescribed fixed cardinality reserve, independent of
$MK$ and compatible volume.

For $W=e^{iw}$, set
\[
                    \Xi_w(y)=\frac1i\log(e^{iy}e^{iw}).
\]
This is a per-bond analytic map with uniformly bounded
derivatives on a fixed small joint ball. With
$E_w=D_y\Xi_w(0)$ and $U_w=D_y^2\Xi_w(0)$,
\begin{equation}
 q_{\rm eff}(w)[u,v]
   =D^2S_{MK}(w)[E_wu,E_wv]
                 +DS_{MK}(w)[U_w(u,v)].
 \label{dgb-coarse-BCH}
\end{equation}
The background-current term is retained.
Compose each scalar support coefficient with $\Xi_w(y)$,
differentiate twice in $y$, and up to twice in $w$ using
fixed spare Cauchy radii. No support enlargement occurs.
Rows, columns and background source slots stay in the
same carrier. A fixed cardinality reserve absorbs their
finite counting factors.

Consequently, for some fixed $\mu,\rho>0$,
\begin{equation}
       \sup_{\|w\|_\infty<\rho}
                    \|q_{\rm eff}(w)\|_\mu\leq R_\mu,
 \label{dgb-coarse-Schur}
\end{equation}
where $\|\cdot\|_\mu$ is the maximum of the exponentially
weighted absolute block row and column sums.
The bound is uniform in $MK$ and volume. After reducing
the exponent to $m>0$, its source derivatives satisfy
\begin{align}
 \sum_{ij}e^{m[d(i,x)+d(j,x)]}
       |(\partial_xq_{\rm eff})_{ij}|&\leq C_1,\notag\\
 \sum_{ij}|(\partial_x\partial_yq_{\rm eff})_{ij}|
                              &\leq C_2e^{-md(x,y)}.
 \label{dgb-coarse-source-bounds}
\end{align}
To obtain these bounds, split the carrier-tree exponential
between its source separation and its row-to-source
distance, then sum that distance. The coarse cell fiber
has fixed dimension, with no factor from $K^4$ fine
coordinates.

\begin{proposition}[Bounded actual coarse determinant contribution]
\label{dgb-coarse-proposition}
Choose a fixed $a>2R_\mu$ and the completion $B=aI$.
The second-background-derivative kernel of
\[
                \tfrac12\mathcal N_0
                        \log\det(q_{\rm eff}+aI)
\]
has a uniform exponential row bound, independent of the
preceding depth, endpoint and compatible volume.
Its true infinite Maxwell coefficient, and its actual
finite reference-normalized coefficients above the
reference threshold, have a uniform absolute bound.
\end{proposition}

\begin{proof}
In the weighted Schur algebra, the Neumann series gives
\begin{equation}
 P_c=q_{\rm eff}+aI,\qquad
                         \|P_c^{-1}\|_\mu\leq 2/a.
 \label{dgb-coarse-inverse}
\end{equation}
The corresponding fine completion is positive on the
real germ: write $x=\xi+H_by$, $C\xi=0$.
Equation \eqref{dgb-effective-Schur} makes the cross term
zero, leaving the positive constrained-divergence form
in $\xi$ and $y^T(q_{\rm eff}+aI)y$.
No uniform fine spectral gap is inferred.

At finite volume the series
\begin{equation}
 \log\det P_c
   =(\dim P_c)\log a+
      \sum_{n\geq1}\frac{(-1)^{n+1}}n
                \operatorname{Tr}(a^{-1}q_{\rm eff})^n
 \label{dgb-coarse-log-series}
\end{equation}
converges absolutely. Let $K_c$ be the Hessian kernel of
$\mathcal N_0\log\det P_c$; the desired half-logarithm
has Hessian $K_c/2$. Explicitly,
\begin{equation}
 K_c(x,y)=\operatorname{Tr}\left[
       P_c^{-1}q_{{\rm eff},xy}
       -P_c^{-1}q_{{\rm eff},x}
                   P_c^{-1}q_{{\rm eff},y}\right].
 \label{dgb-coarse-trace}
\end{equation}
The first trace is bounded by the second line of
\eqref{dgb-coarse-source-bounds} and the inverse bound.
For the second trace, use the two source-rooted first
derivative bounds and spend a spare inverse-edge
exponent to connect their roots. Absolute summation gives
\begin{equation}
                            |K_c(x,y)|
                                \leq C_a e^{-m'd(x,y)}
 \label{dgb-coarse-Hessian-decay}
\end{equation}
for a fixed $m'>0$, uniformly in depth and volume.
Periodic sums are bounded by their infinite-lattice
counterparts. Its second displacement moment is therefore
uniformly bounded. The finite reference denominators
are uniformly bounded below above their established
threshold, which proves the finite statement as well.
\end{proof}

The actual action's local gauge identity makes its
relative Hessian transform by orthogonal bond-frame
conjugation under real local gauges. Thus the coarse
massive determinant has the appropriate total local
gauge identity on its common germ; no coefficientwise
Ward identity was used in the proof.
With $B=aI$, \eqref{dgb-effective-determinant} now reduces
the unbounded-endpoint comparison to the fine combination
\[
 \log\det G_b-\tfrac12\log\det\!
       \left[s''-\lambda\cdot f''
                     +D_bD_b^T+aC^TC\right],
\]
up to the proved uniformly bounded coarse contribution.
For a general completion the full-site vector difference
is $-\lambda\cdot f''+C^TBC-\eta^2I$. For this fixed-$a$
choice it is
\begin{equation}
              Q_f-P_\eta
                    =-\lambda\cdot f''+aC^TC-\eta^2I.
 \label{dgb-fixed-a-difference}
\end{equation}
Equation \eqref{dgb-reference-difference} instead records
the canonical completion $B=I-q_{\rm eff}$.
Each comparison uses all factors from its one chosen
completion.

\subsection{The remaining uniform comparison}
\label{dgb-scope-section}

The deposited outline asserts a bounded error in the
Maxwell coefficient when its composite determinant is
compared with a massive vector/ghost reference, at
$\eta=\ell^{-k}$. The exact identities above do not
establish that error bound. In particular, the formal
background effective-action convention and massless
nonzero-external-momentum amplitudes of a reference
calculation do not identify this constrained normalizer
automatically.

For a precisely matched reference, after removing the
bounded coarse effective-Hessian contribution, the
full-site formulation must still control the Maxwell
moment of the coarse-site package
\[
                         -\log\det S_g+\tfrac12\log\det S_v,
\]
together with the global-base and ghost-mass correction,
and the Maxwell moment of
$\log\det Q_f-\log\det P_\eta$ with all multiplier,
coarse-completion and mass terms for the selected
completion. They must also identify
the reference background with the actual smooth pre-gauge
minimum and preserve the actual finite extraction and
zero-mode conventions. Based-gauge covariance allows
evaluating \eqref{dgb-main-identity} at that pre-gauge
representative, but does not replace it by an arbitrary
external interpolation.

Each determinant Hessian has the concrete formula
\begin{equation}
 D^2\log\det P[u,v]
   =\operatorname{Tr}\left(
       P^{-1}P_{uv}-P^{-1}P_uP^{-1}P_v\right).
 \label{dgb-trace-interface}
\end{equation}
The established fixed-$K$ bounds make the relevant
second-displacement-moment trace sums meaningful. Their
constants may depend on $K$; they do not prove an $O(1)$
bound as $K\to\infty$.
The soft mode \eqref{dgb-Rayleigh} shows why a separate
uniform inverse estimate cannot supply this step.
A coupled cancellation or a scale-resolved comparison
is still required. No universal coefficient, sign,
averaged drift or ultraviolet coupling trajectory is
deduced here.


% No preamble. Prefix olc-. Fixed-regulator leading stationary amplitudes.
\section{The actual leading determinant correction satisfies a cocycle}
\label{olc-section}

This section proves an exact identity of normalized leading stationary
amplitudes at a fixed finite regulator. It compares these amplitudes
through one auxiliary local orbit-tube integral. It does not identify
direct and sequential sharp-window sectors at positive coupling.
All maps, pivot factors and retained Hessians are those of the
actual composite ordinary constraint.

\subsection{Objects and statement}
\label{olc-setup-section}

Fix a finite compatible periodic regulator, a preceding homogeneous
depth $M$, and two allowed ratios $R=\ell^r$, $L=\ell^s$.
Let $X,Y,Z$ be the current, $R$-coarse and $RL$-coarse link
configuration spaces, each with product normalized Haar measure.
The actual iterated ordinary maps satisfy
\begin{equation}
 f=F_R:X\longrightarrow Y,\qquad
 h=F_L:Y\longrightarrow Z,\qquad
 h\circ f=F_{RL}.
 \label{olc-maps}
\end{equation}
No equality with a direct-$R$ or direct-$RL$ ordinary average
is used.

The coupling-free action $S_M$ is the established homogeneous
ordinary minimum action, normalized at the identity. For an
allowed ratio $K$, write $V_*^{M,K}(W)$ for the physical
gauge-zero stationary representative, using the independently
chosen direct-$K$ internal contour gauge $P_K$. Let
$A_{M,K}(W)$ be the full retained Hessian of $S_M+G_K$ there,
in the actual constraint chart. Let $d_K(V_*)$ be the product
of inverse pivot Jacobians. Define
\begin{equation}
 D_{M,K}(W)
  =\log\frac{d_K(V_*^{M,K}(W))}{d_K(1)}
   -\frac12\log\det\!\left[
         A_{M,K}(W)A_{M,K}(0)^{-1}\right].
 \label{olc-leading}
\end{equation}
Here $D_{M,K}(1)=0$. These are the actual $\Delta+\mathcal Q_{\det}$
terms, with positive real branches near the identity and their
given analytic continuations. Constraint curvature and the
absolute contour-gauge Hessian are included in $A_{M,K}$.

Let $V_p^{MR,L}(W)$ be the outer pre-root-gauge stationary
representative, and let $\mathcal T_{MR,L}$ denote substitution
into this field.

\begin{theorem}[Finite leading-amplitude cocycle]
\label{olc-theorem}
On a real small-field germ,
\begin{equation}
 D_{M,RL}
   =\mathcal T_{MR,L}D_{M,R}+D_{MR,L}.
 \label{olc-cocycle}
\end{equation}
The identity extends to every connected common analytic domain
already available for its two sides. The germ used in its proof
may depend on the fixed regulator and the ratios.
\end{theorem}

The proof keeps all Haar, gauge and Gaussian normalization
factors. They disappear only when each leading amplitude is
divided by its value at the identity.

\subsection{Based gauge groups and an exact insertion}
\label{olc-groups-section}

Let $H_R$ consist of current site gauges equal to the identity
at every intermediate $R$-center, and let $H_{RL}$ be based
at the final $RL$-centers. Restriction to intermediate centers
gives
\begin{equation}
 1\longrightarrow H_R\longrightarrow H_{RL}
       \longrightarrow H_L\longrightarrow1,
 \label{olc-group-sequence}
\end{equation}
where $H_L$ is the intermediate group based at final centers.
The site coordinates split into the non-intermediate sites
and the intermediate non-final centers. Embed $H_L$ by setting
its other current-site values equal to the identity. The two
factors have disjoint supports, so their direct product is
$H_{RL}$ and normalized Haar is their product Haar.
There is no background-dependent orbit-volume factor.

Each based action is free. A gauge preserving every bond of a
connected graph is determined by its value at one site; its
value at a based center is the identity.
The actual maps obey
\begin{equation}
 f(U^u)=f(U)^{u|_{\text{intermediate centers}}},\qquad
 h(V^v)=h(V)^{v|_{\text{final centers}}}.
 \label{olc-equivariance}
\end{equation}
Thus $f$ is $H_R$-invariant, $h$ is $H_L$-invariant, and
$h\circ f$ is $H_{RL}$-invariant. Product link Haar is
invariant under all these site actions.

For $u\in H_K$, endpoint covariance of the actual direct-$K$
contour variables gives
\[
 P_K(U^u;y,x)=P_K(U;y,x)u(x)^{-1}.
\]
With the actual nonnegative class penalty
$q(v)=1-\operatorname{Re}\tau(v)$, whose unique zero is the
identity, put $G_K(U)=\sum_{x\ {\rm not\ a\ center}}q(P_K(U;y,x))$.
Then
\begin{equation}
 \int_{H_K}e^{-G_K(U^u)/g^2}\,du
       =c(g)^{n_{s,K}},\qquad
 c(g)=\int_G e^{-q(v)/g^2}\,dv .
 \label{olc-gauge-insertion}
\end{equation}
It is therefore exact to insert $e^{-G_K/g^2}$ into an
$H_K$-invariant tube integral and multiply by
$c(g)^{-n_{s,K}}$. This factor is independent of the target
field. Its asymptotic constant and power of $g$ need not be
computed below.

\subsection{A consistent local orbit-tube phase}
\label{olc-tube-section}

The original $S_M$ is known on a small-link domain. We first
justify extending it to the particular saturated tube needed
for an exact pushforward composition.

Suppose two sufficiently small configurations $U,V$ satisfy
$V=U^u$ with $u\in H_{RL}$. The bond equation, in a
unitary-invariant matrix norm, yields
\begin{equation}
 \|u(x)-u(x+e_\mu)\|
   \leq C\bigl(\|U(x,\mu)-I\|+\|V(x,\mu)-I\|\bigr).
 \label{olc-gauge-edge}
\end{equation}
Every site has a path of at most $C_dRL$ edges to a final
center, where $u=I$. Consequently
\begin{equation}
 \|u-I\|_\infty
   \leq C_dRL\bigl(\|U-I\|_\infty+\|V-I\|_\infty\bigr).
 \label{olc-gauge-small}
\end{equation}
Shrink the comparison ball so that the near-identity logarithm
$\lambda=\log u$ exists. Its edge differences satisfy the
same type of estimate as \eqref{olc-gauge-edge}.
Along $u_t=\exp(t\lambda)$ the transformed links stay within
a fixed multiple of the endpoint smallness. Fixed color-norm
conversions are included in this constant. Shrinking the ball
again keeps the whole path in the established action domain.
Local Ward invariance therefore gives
\begin{equation}
                         S_M(U)=S_M(V).
 \label{olc-phase-consistency}
\end{equation}
The same argument applies on the intermediate graph and to
$H_L$. No radius uniform in $RL$ is claimed in this step.

On the $H_{RL}$-saturation of this small representative
neighborhood, define
\begin{equation}
                       S_{\rm ext}(U^u)=S_M(U).
 \label{olc-extended-phase}
\end{equation}
Equation \eqref{olc-phase-consistency} makes the definition
independent of the small representative. The free compact
group action has an actual local slice $P_{RL}=I$: the unique
based transformation to it is $u(x)=P_{RL}(U;y,x)$.
Hence $S_{\rm ext}$ is smooth, and real analytic in local
slice coordinates, and is invariant on the saturation.

The local ordinary maps and contour variables extend
equivariantly to this tube. One may use their endpoint-covariant
relative-log formulas or define them from the slice. On
overlaps the two definitions agree by
\eqref{olc-equivariance} and the original local formulas.
No measurable completion away from the tube is invoked.

The proved physical branch is a strict local constrained
minimum modulo the based orbit. Its actual retained
Hessian after adding $G_K$ is positive. On the slice $P_K=I$,
the gauge penalty is identically zero, so that positivity
also gives positivity of the action Hessian on the quotient
constrained tangent. The finite-dimensional implicit-function
argument yields a unique local critical orbit in a smaller
neighborhood.

Choose a compact closure of a sufficiently small real target
germ. Its stationary orbits form a compact set. Choose a smooth
$H_{RL}$-invariant cutoff $\chi$, equal to one near these orbits
and supported in a tube containing no other quotient saddle.
It can be chosen one near every inner conditional minimum
over a smaller compact outer saddle neighborhood as well:
use continuity of the minimum maps, choose nested neighborhoods,
and average the cutoff over the compact based action.
The averaging leaves it one on the invariant inner neighborhood.
The quotient phase has a positive gap on the remaining compact
part of the support.

This cutoff is an auxiliary device for determining a leading
coefficient. It is neither the actual retained sharp window
nor a proposed global Wilson density.

\subsection{Exact Haar pushforward composition}
\label{olc-pushforward-section}

For $g>0$, consider the positive finite measure
\begin{equation}
              \chi(U)e^{-S_{\rm ext}(U)/g^2}\,d\mu_X(U).
 \label{olc-tube-measure}
\end{equation}
The maps $f,h,h\circ f$ are submersions near all relevant
orbits. At a small representative, the actual pivot blocks
give full rank; equivariance preserves rank on the orbit.
A smaller tube retains the property. Its local pushforward
densities are thus smooth relative to product Haar.

Write $J_g(V)$ for the pushforward density through $f$, and
$I_g(W)$ for the density through $h\circ f$. Exact
pushforward composition gives
\begin{equation}
                I_g(W)=\int_{h(V)=W}J_g(V)\,d\mu_h(V).
 \label{olc-Fubini}
\end{equation}
Here $d\mu_h$ is the conditional coarea density relative to
product Haar. Equivalently, test both sides against a smooth
function of $W$ and use ordinary Fubini. No multiplication
of formal delta distributions is required.

The first pushforward is exactly $H_L$-invariant. Indeed,
lift a based intermediate gauge by
\eqref{olc-group-sequence}, use the $H_{RL}$ invariance of
$\chi$ and $S_{\rm ext}$, and change the current Haar variables.
This is precisely the invariance needed in the outer integral.

The actual two-minimum comparisons identify the inner
critical value and then the outer one:
\begin{equation}
 \min_{f(U)=V}S_{\rm ext}(U)=S_{MR}(V),\qquad
 \min_{h(V)=W}S_{MR}(V)=S_{MRL}(W).
 \label{olc-minimum-values}
\end{equation}
These are local comparisons on the chosen tube.
The conditional minima depend smoothly on the outer
parameters modulo the inner based orbit. The cutoff is one
there. Outside the smaller outer saddle neighborhood, the
compact quotient phase gap makes the contribution
exponentially smaller. All uniformity here concerns the
fixed finite regulator and fixed ratios only.

\subsection{The finite-dimensional Laplace fact}
\label{olc-Laplace-section}

We record the elementary asymptotic used twice. In a finite
real chart suppose
\[
 F(z,w)=s(w)+\tfrac12z^TA(w)z+O(|z|^3),\qquad A(w)>0,
\]
and the critical point is the unique local minimum on the
chosen compact support. The density $m(z,w)$ is smooth and
its cutoff is one near zero. On a small compact parameter
germ the Hessian has a common positive lower bound. Then
\begin{equation}
 \int e^{-F(z,w)/g^2}m(z,w)\,dz
  =e^{-s(w)/g^2}g^n(2\pi)^{n/2}
        m(0,w)\det A(w)^{-1/2}[1+o(1)]
 \label{olc-Laplace}
\end{equation}
uniformly in that parameter germ.

To prove it, shrink the chart until $F-s\geq c|z|^2$.
Set $z=g\xi$, first restrict $\xi$ to a fixed bounded set,
and use Taylor's formula for uniform convergence to the
Gaussian. A common Gaussian majorant removes that restriction.
The compact remainder outside a smaller chart has a positive
phase gap and is exponentially negligible. This proof also
allows a smooth positive amplitude, whose leading value is
its value at the saddle, and an amplitude family converging
uniformly there with a bounded relative error.

For an orbit integral, first insert
\eqref{olc-gauge-insertion}. The phase becomes $S+G_K$.
Its unique saddle is the gauge-zero representative: the
quotient phase has its strict minimum there modulo the orbit,
and on the compact based orbit $G_K$ has its unique zero
at $P_K=I$. Away from a neighborhood of that representative,
the sum of the quotient and gauge penalties has a positive
gap. The actual retained chart therefore reduces the leading
asymptotic to \eqref{olc-Laplace}, with the exact factor
$c(g)^{-n_{s,K}}$ still present. No unspecified quotient
volume form or ultraviolet-uniform error estimate is used.

\subsection{Identification with the actual pivot and determinant}
\label{olc-identification-section}

The exact local constrained measure has, at its saddle,
the factor
\begin{equation}
                       h_0^{\,n_e-n_c}d_K(V_*).
 \label{olc-saddle-measure}
\end{equation}
Indeed relative Haar coordinates are
$d(e^{iB})=h_0j(B)\,dB$ with $j(0)=1$; the normalized
fluctuation straightening Jacobian is one at zero; and
solving the constraint contributes the inverse pivot product.
The target density is relative to group Haar. Consequently
there is no extra factor $j(\log W)$ in
\eqref{olc-saddle-measure}.

Let $I_{M,K;g}$ denote an appropriate local tube pushforward.
The phase-corrected, identity-normalized stationary amplitude is
\begin{equation}
 e^{D_{M,K}(W)}
   =\lim_{g\downarrow0}
       e^{[S_{MK}(W)-S_{MK}(1)]/g^2}
          \frac{I_{M,K;g}(W)}{I_{M,K;g}(1)} .
 \label{olc-normalized-amplitude}
\end{equation}
The gauge integral, powers of $g$, Haar constants, flat
determinants and all scalar cutoff factors cancel in this
ratio. The remaining saddle measure and Hessian are exactly
\eqref{olc-leading}. A cutoff equal to one near the critical
orbit changes no leading coefficient.

The actual sharp-window integral has the same saddle,
pivot factor and retained Hessian. Its variable is $z=g\xi$
and its window is the actual product window of threshold
$T(g)$. Choose $T(g)\to\infty$ with $gT(g)^3\to0$.
The established centered Gaussian estimate makes its
normalized logarithmic expectation
\begin{equation}
                          O\!\left(\tau(T(g))+g^2T(g)^6\right)
                              \longrightarrow0.
 \label{olc-sharp-leading}
\end{equation}
Thus the sharp-window and the auxiliary fixed-tube integrals
have the same normalized leading amplitude. This does not
equate the integrals at positive $g$.

The auxiliary first pushforward supplies only the intermediate
based $H_L$ invariance: its input was made $H_{RL}$-invariant.
The phase and normalized leading amplitude therefore have
that identity on the needed intermediate orbit. It suffices
to replace the rooted outer saddle by $V_p$.
The cutoff has not been made invariant under every output
site gauge with nontrivial final-center values. The full
local output gauge identity for the actual $D$ is the
previously proved one from the zero-input quantum leading
limit, not an additional invariance claimed for this cutoff.

\subsection{Sequential saddle evaluation}
\label{olc-proof-section}

\begin{proof}[Proof of Theorem~\ref{olc-theorem}]
Apply \eqref{olc-Laplace} on the $f$ fibers, uniformly over
a compact neighborhood of the outer stationary orbit.
With a scalar $c_1(g)$ independent of $V$,
\begin{equation}
 J_g(V)=c_1(g)e^{-S_{MR}(V)/g^2}
                   e^{D_{M,R}(V)}[1+o(1)] .
 \label{olc-inner-asymptotic}
\end{equation}
The relative error is uniform. Exact $H_L$ invariance reduces
the full compact outer orbit to a small slice. The inner
cutoff is one at every relevant conditional minimum.

Insert \eqref{olc-inner-asymptotic} into
\eqref{olc-Fubini}. The outer Laplace formula evaluates the
smooth amplitude $e^{D_{M,R}(V)}$ at its existing classical
minimum. With another scalar $c_2(g)$,
\begin{equation}
 I_g(W)=c_1(g)c_2(g)e^{-S_{MRL}(W)/g^2}
       e^{D_{M,R}(V_p^{MR,L}(W))+D_{MR,L}(W)}[1+o(1)] .
 \label{olc-sequential-asymptotic}
\end{equation}
The amplitude gradient has no factor $g^{-2}$ and hence
does not alter the leading phase saddle. The outer rooted
representative gives the same amplitude by the based
$H_L$ invariance. The compact phase gap controls the
remaining part of the outer integration.

Direct evaluation of the same integral through $h\circ f$
instead gives a scalar $c_{12}(g)$ and
\begin{equation}
 I_g(W)=c_{12}(g)e^{-S_{MRL}(W)/g^2}
                         e^{D_{M,RL}(W)}[1+o(1)] .
 \label{olc-direct-asymptotic}
\end{equation}
Normalize both equations by their values at $W=I$ after
removing the phase. There $V_p=I$ and every $D$ is zero.
All scalar factors cancel, without any computation of the
relation between $c_{12}$ and $c_1c_2$. Taking $g$ to zero
yields
\[
 e^{D_{M,RL}(W)}
   =e^{D_{M,R}(V_p^{MR,L}(W))+D_{MR,L}(W)} .
\]
On the positive real germ the logarithms are real and
zero at the identity. They therefore satisfy
\eqref{olc-cocycle}, without a branch ambiguity.
Both sides are given holomorphic pivot/determinant
functionals. The identity theorem extends their germ
identity to each connected common analytic domain.
The smaller comparison germ used by Laplace does not
shrink a domain already established for those functions.
\end{proof}

\subsection{Joint infinite systems and the Maxwell coefficient}
\label{olc-limit-section}

For fixed $R,L$, the existing modules give uniform anchored
analytic families for all three $D$ terms. On a joint local
limit system the finite identity passes on compact fields
by local normal convergence and anchored tail bounds;
the same bounds justify the classical substitutions.
It gives a scalar quotient cocycle, not a coefficientwise
identity of carrier decompositions.

Apply exact preservation of the true infinite Maxwell
coefficient under classical transport. Then
\begin{equation}
 \beta_{MRL}(D_{M,RL})
   =\beta_{MR}(D_{M,R})+\beta_{MRL}(D_{MR,L}).
 \label{olc-Maxwell-cocycle}
\end{equation}
Thus a long fixed-ratio leading coefficient is the sum of
its consecutive step coefficients in this joint system.
This proves neither their sign nor their convergence as
the preceding depth grows.

At finite volume, the symmetric short-cutoff moment
$\beta_{M,N}$ does not commute exactly with classical
transport. One must retain the separately controlled
finite leakage when applying it to
\eqref{olc-cocycle}. No exact finite scalar
$\beta$ additivity is asserted.

\subsection{Scope}
\label{olc-scope-section}

The result rests on a consistent local orbit-tube extension,
exact finite Haar pushforward composition, and ordinary
finite-dimensional saddle asymptotics using the actual
pivot, contour penalty and retained Hessian. It assumes
no ultraviolet-uniform stationary-phase remainder.

It does not identify direct and sequential sharp-window
sectors at positive $g$, extend the original local action
to every gauge configuration, or identify every local
sector with the full Wilson measure. The auxiliary cutoff
disappears only from the leading coefficient.

The background-field/ghost determinant comparison, a
bounded comparison error in depth, the numerical
logarithmic coefficient, its sign, physical coupling
tuning and the treatment of the other field sectors
remain separate from this cocycle.


% No preamble. Prefix fbt-. Conditional qualitative positive block transfer.
\section{Qualitative transfer of a positive leading block drift}
\label{fbt-section}

The compactness and moment-passage inputs in this section are
available for the actual homogeneous local macrostep. Its
positivity hypothesis is explicitly additional and unproved.
The conclusion requires no quantitative finite-volume rate.

\subsection{The actual fixed-length coefficients}
\label{fbt-setup}

Fix the primitive ratio, elementary composite ratio \(R\),
grouping \(J\), common analytic domains and extraction
convention. Fix a finite \(q\ge1\), chosen to be a multiple
of the activity-envelope block length \(q_a\) when applying
the tuning consequence below. Let \(\mathcal P\) be the
compatible nested cubic period family: every fixed power
of \(R\) divides all sufficiently late periods in
\(\mathcal P\).

For a \(q\)-step block with final period \(N\in\mathcal P\),
put
\begin{equation}
 M_j=MR^j,\qquad Q_j=R^{q-j}N,\qquad 0\le j\le q.
 \label{fbt-block-scales}
\end{equation}
Here \(M\) is any permitted preceding homogeneous depth.
Let \(D_{M_j,R;Q_{j+1}}\) be the actual finite leading
correction \(\Delta+\mathcal Q_{\det}\) of step \(j\).
Define
\begin{equation}
 b_{0,j}^{M,N}
    =\beta_{M_jR,Q_{j+1}}(D_{M_j,R;Q_{j+1}}),
 \qquad
 B_q(M,N)=\sum_{j=0}^{q-1}b_{0,j}^{M,N}.
 \label{fbt-finite-coefficients}
\end{equation}
Each functional is the actual finite short-lift ratio
normalized by its actual finite classical reference.
These coefficients are real on the real branches.
No finite ratio is replaced by an infinite moment.

The uniform outer-domain leading-family bound and the
finite moment estimate give, above the fixed reference
threshold,
\begin{equation}
                          |b_{0,j}^{M,N}|\le B_0
 \label{fbt-leading-bound}
\end{equation}
uniformly in depth and permitted volume. The reference
denominators are bounded below by \(1/2\).

\subsection{Varying-depth structural limits and their moments}
\label{fbt-compactness}

Take any \(M_n\) and any \(N_n\in\mathcal P\) tending to
infinity. The old depths need not be bounded or converge.
Extract jointly the classical references and background
maps at the shifted depths \(M_nR^j\), their required
fixed-order vertices, the actual retained flat precisions
and covariances, and all \(q\) leading corrections in
\eqref{fbt-finite-coefficients}.
The uniform strict support, common-domain and exponential
source-Hessian bounds provide such a subsequence. If a
structural reference comparison requires additional fixed
depth shifts or endpoint powers, include their countable
collection in the same diagonal extraction. Cofinal
divisibility guarantees their eventual finite availability.

Write the limiting corrections as \(D_j^*\), with true
infinite Maxwell functional \(\beta_{j+1}^*\) at output
scale \(j+1\), and set
\begin{equation}
                       B_q^*=\sum_{j=0}^{q-1}
                                         \beta_{j+1}^*(D_j^*).
 \label{fbt-limit-sum}
\end{equation}
The scale indices describe the structural limit. They
do not identify it with one fixed finite-depth system.

These are actual determinant structural limits. At fixed
\(R\) the retained fiber per current block has fixed finite dimension;
the whole retained operator still has dimension growing with volume.
Local matrix convergence and exponential envelopes pass
products and the inverse identity \(A_0C_0=I\) by truncating
the intermediate sums. The chart and pivot maps are fixed;
common-domain Cauchy estimates give convergence of classical
vertices through order four. The absolutely convergent
trace formula therefore identifies the limiting Hessian
with the same actual determinant expression. One may
equivalently extract the strict scalar families jointly
and identify their Hessians in this way. No old fine
coordinate count occurs.

\begin{lemma}[Passage of the actual finite moments]
\label{fbt-moment-lemma}
Along every such joint subsequence,
\begin{equation}
 b_{0,j}^{M_n,N_n}\longrightarrow
                      \beta_{j+1}^*(D_j^*)\quad(0\le j<q),
 \qquad
 B_q(M_n,N_n)\longrightarrow B_q^*.
 \label{fbt-moment-limits}
\end{equation}
\end{lemma}
\begin{proof}
For a fixed \(j\), the leading source-Hessian kernels
converge locally and have a common exponential row
bound. Restrict their short-lift second moment first
to a fixed displacement radius. The finite sum converges.
The rest has a uniformly summable exponential
second-moment tail. The finite short-lift cutoff tends
to infinity since \(Q_{j+1}=R^{q-j-1}N_n\to\infty\).

The actual finite reference denominator tends uniformly
to one by the flat normalized classical periodization
estimate, even for these varying depths. Division proves
the first assertion. The sum has fixed finite length
\(q\), proving the second. Neither an extensive total
limit nor a regulator convergence rate was used.
\end{proof}

\subsection{The conditional positive transfer}
\label{fbt-positive-section}

\paragraph{Unproved structural positivity hypothesis.}
Assume there is one \(\delta>0\) such that \emph{every}
joint structural infinite limit obtainable above satisfies
\begin{equation}
                                  B_q^*\ge\delta.
 \label{fbt-positive-hypothesis}
\end{equation}
The same margin must work for all varying-depth and
permitted-period sequences. Positivity in one selected
fixed-depth limit is insufficient.

\begin{proposition}[Uniform finite block lower bound]
\label{fbt-transfer-proposition}
Under \eqref{fbt-positive-hypothesis}, for every
\(0<\varepsilon<\delta\) there is a threshold
\(N_*(\varepsilon)\), independent of preceding depth,
such that
\begin{equation}
 B_q(M,N)\ge\delta-\varepsilon
 \quad\hbox{for all permitted }M,\quad
 N\in\mathcal P,\quad N\ge N_*(\varepsilon).
 \label{fbt-finite-lower}
\end{equation}
In particular the finite block is at least \(\delta/2\)
above one uniform period threshold.
\end{proposition}
\begin{proof}
Otherwise choose \(N_n\to\infty\) in \(\mathcal P\)
and old depths \(M_n\) such that
\(B_q(M_n,N_n)<\delta-\varepsilon\).
Take a joint structural subsequence. By
Lemma~\ref{fbt-moment-lemma}, its limiting sum is at most
\(\delta-\varepsilon\), contradicting
\eqref{fbt-positive-hypothesis}.
\end{proof}

The threshold is also uniform in a block's starting
stage: absorb that stage into its preceding depth \(M\).
Its final period is the \(N\) in
\eqref{fbt-block-scales}. The proof gives no numerical
threshold or convergence rate.

\subsection{The finite tuning input needs only this qualitative bound}
\label{fbt-tuning-section}

Take \(q\) to be a multiple of \(q_a\). Write
\(\mathcal F_{j,g}\) for the actual finite quantum extraction map,
\(x_j\) for the actual real extracted input class, and
\(d_j(g,x)=\beta_{j+1}(\mathcal F_{j,g}(x))\) for its
inverse-coupling decrement. On the established
finite activity envelope the actual decrement obeys
\begin{align}
 |d_j(g,x_j)-b_{0,j}|
 \le{}&C_{\rm int}\omega(N_{\min})
       +C_{{\rm par},\beta}
                    \{\tau(T(g))+g^2T(g)^6\}\notag\\
 &+C_{{\rm in},\beta}\,gT(g)/b.
 \label{fbt-actual-error}
\end{align}
The first term tends qualitatively to zero as the
period threshold increases. The other terms tend to
zero in the weak-coupling limit at the fixed field
radius and chosen real threshold \(T(g)\).

Use \eqref{fbt-finite-lower} with
\(\varepsilon=\delta/2\). Raise the period threshold
and lower the weak-coupling ceiling until
\eqref{fbt-actual-error} is at most \(\delta/(4q)\).
Then every permitted \(q\)-step block satisfies
\begin{equation}
 \frac{\delta}{4}\le
       \sum_{j=k}^{k+q-1}d_j(g_j,x_j)
       \le qD_{\rm abs}.
 \label{fbt-actual-block}
\end{equation}
The upper bound is the proved absolute-increment bound.
Repeating the \(q_a\)-step activity blocks gives the same
intermediate envelope throughout this longer block.
No new large-endpoint field-dilation estimate is used.

These are the positive block bounds used in the
continuous clipped-recursion shooting argument.
Put \(\Gamma_{\rm w}=g_{\rm w}^{-2}\), where
\(g_{\rm w}\) is the chosen weak-coupling ceiling.
For a terminal inverse coupling
\(\Gamma_*\ge\Gamma_{\rm w}+qD_{\rm abs}\), that argument
then gives an initial value in
\begin{equation}
 [\,\Gamma_*+m\delta/4,\ \Gamma_*+mqD_{\rm abs}\,]
 \label{fbt-shooting-interval}
\end{equation}
which makes \(mq\) genuine weak local extraction steps
and reaches \(\Gamma_*\). This remains conditional on
\eqref{fbt-positive-hypothesis}.

If an averaged formulation for arbitrary finite
interval lengths is desired, write \(n=mq+r\),
\(0\le r<q\), and use \eqref{fbt-leading-bound} on
the final fringe. For intervals entirely above the
common period threshold,
\begin{equation}
 \sum_{j=k}^{k+n-1}b_{0,j}
 \ge cn-C_{\rm fin},\qquad
 c=\frac{\delta}{2q},\qquad
 C_{\rm fin}=(q-1)(B_0+c).
 \label{fbt-optional-fringe}
\end{equation}
Indeed the complete blocks contribute at least
\(m\delta/2\), and the fringe contributes at least
\(-rB_0\). The shooting proof needs only the fixed-block
bound, so this additional formulation is optional.

\subsection{The remaining sign problem}
\label{fbt-scope}

Include the one long-ratio \(R^q\) determinant in the
joint extraction. The proved leading cocycle and true
infinite classical marginal preservation give
\begin{equation}
                         B_q^*
                          =\beta_q^*(D_{0,R^q}^*).
 \label{fbt-long-ratio}
\end{equation}
A uniformly positive reference comparison for this
coefficient in every structural limit would therefore
supply \eqref{fbt-positive-hypothesis}. The equality
itself proves no sign.

Once that positive structural margin is established,
qualitative compactness is enough for the needed finite
block lower bound; a separate exponential regulator
rate is unnecessary. This section does not establish
the margin, exact finite marginal additivity, summability
of errors over an infinite history, uniqueness of the
old-depth limit, or convergence of tuned local densities
to a continuum measure.


% No preamble. Prefix fts-. Conditional finite endpoint tuning only.
\section{Finite ultraviolet tuning under an averaged drift hypothesis}
\label{fts-section}

The finite activity induction, leakage estimate, zero-input
parity remainder and real-coupling continuity are available.
This section gives a conditional endpoint argument for their
actual finite local maps. Its additional hypothesis is a
uniformly positive averaged drift of the \emph{actual finite}
leading determinant coefficients. That hypothesis is not
proved here.

\subsection{The actual decrement on the induction envelope}
\label{fts-decrement-section}

Fix the finite-induction data: elementary ratio $R$, grouping
$J$, radius $b$, carrier exponent $a_0$, amplitude block length
$q_a$, endpoint radius $r_b=K_{\mathrm{amp}}b^2$, and intermediate
bound $C_{\mathrm{int}}b^2$. All references and $\beta_j$ are the
actual finite reference-normalized ones. For a real extracted
input define
\begin{align}
 d_j(g,x)&=\beta_{j+1}(F_{j,g}(x)),\notag\\
 \gamma_{j+1}&=\gamma_j-d_j(g_j,x_j),\qquad
 g_j=\gamma_j^{-1/2},\notag\\
 x_{j+1}&=P_{j+1}F_{j,g_j}(x_j).
 \label{fts-actual-update}
\end{align}
Real branches and real initial data give real decrements.
The activity induction applies to every realized weak
coupling sequence while the permitted periods stay above
its threshold. In particular $x_0=0$ is allowed.

Write
\begin{equation}
 b_{0,j}=\beta_{j+1}
       \bigl(D_{M_j,R;N_{j+1}}\bigr)
 \label{fts-leading-coefficient}
\end{equation}
for the finite pivot/Gaussian leading correction. It is
independent of $g$ and of the inherited interaction, but
depends on the actual classical data and regulator. It is
not replaced by an independently selected infinite coefficient.

Since every $x_j$ is extracted, the finite leakage bound,
parity estimate and quantum input comparison give
\begin{align}
 |d_j(g,x_j)-b_{0,j}|
 \leq{}& C_{\mathrm{int}}\omega(N_{\min})
   +C_{\mathrm{par},\beta}
       \{\tau(T(g))+g^2T(g)^6\}\notag\\
 &+C_{\mathrm{in},\beta}\frac{gT(g)}b .
 \label{fts-decrement-error}
\end{align}
The first term is the finite classical leakage, the second
includes the fixed outer-radius factor, and the third includes
$C_{\mathrm{int}}$ and the beta bound. This follows from
$F=T+G$ and $G(x)=G(0)+(G(x)-G(0))$. The $b^2$ in the
inner-radius remainder cancels beta's $b^{-2}$ bound.
No exact finite identity $\beta_{\mathrm{out}}T
=\beta_{\mathrm{in}}$ is used.

At fixed $b,R,J$, any prescribed tolerance $\epsilon>0$
in \eqref{fts-decrement-error} can be obtained by raising
one period threshold and then lowering a coupling ceiling.
The choices are independent of the number of permitted
steps. The finite-induction magnitude bound also supplies
\begin{equation}
                 |d_j(g,x_j)|\leq D_{\mathrm{abs}}
 \label{fts-absolute-increment}
\end{equation}
uniformly on the envelope and in the allowed history lengths
and regulators.

\subsection{The additional hypothesis and positive block increments}
\label{fts-hypothesis-section}

\paragraph{Unproved averaged-drift hypothesis.}
Assume that there are constants $c>0$ and
$C_{\mathrm{fin}}\geq0$, independent of old depth and the
permitted finite periods, such that
\begin{equation}
       \sum_{j=k}^{k+n-1}b_{0,j}\geq cn-C_{\mathrm{fin}}
 \label{fts-drift-hypothesis}
\end{equation}
on the contemplated stage intervals. It suffices to assume
this for the block length and starting stages used below.
No pointwise positivity of an individual coefficient is required.

Choose $q$ to be a positive multiple of $q_a$ with
$C_{\mathrm{fin}}\leq cq/4$. Strengthen the period and
coupling thresholds until the error in
\eqref{fts-decrement-error} is at most $c/4$. Every $q$-step
sequence of projected maps started in the endpoint ball,
at arbitrary prescribed couplings in that weak range, then
satisfies
\begin{equation}
 B_-\leq\sum_{j=k}^{k+q-1}d_j(g_j,x_j)\leq B_+,
 \qquad
 B_-=\frac{cq}{2}>0,\qquad B_+=qD_{\mathrm{abs}}.
 \label{fts-block-decrements}
\end{equation}
Indeed the lower bound is
$cq-C_{\mathrm{fin}}-q\epsilon\geq cq/2$, and the upper
bound follows from \eqref{fts-absolute-increment}.

Each consecutive $q_a$ subblock returns to the same endpoint
ball. Thus the same $C_{\mathrm{int}}b^2$ envelope controls
the entire larger block. This step repeats already proved
quantum blocks; it does not invoke a new direct large-endpoint
high-degree estimate or require a new field radius depending
on $q$.

\subsection{A continuous auxiliary clipped recursion}
\label{fts-clipping-section}

Let $\Gamma_{\mathrm w}=g_{\mathrm w}^{-2}>0$ be the
strengthened weak threshold. For the shooting proof only,
evaluate the map for every real scalar $\gamma$ at
\begin{equation}
 g_{\mathrm{eval}}(\gamma)
       =\max\{\gamma,\Gamma_{\mathrm w}\}^{-1/2},
 \label{fts-clipped-evaluation}
\end{equation}
while leaving the scalar update unclipped:
\begin{align}
 \gamma_{\mathrm{next}}
       &=\gamma-d_j(g_{\mathrm{eval}}(\gamma),x),\notag\\
 x_{\mathrm{next}}
       &=P_{j+1}F_{j,g_{\mathrm{eval}}(\gamma)}(x).
 \label{fts-clipped-update}
\end{align}
This is auxiliary when $\gamma<\Gamma_{\mathrm w}$; no
physical density interpretation is assigned there.
Its evaluated couplings always belong to the proved weak
range. Starting in the endpoint ball, finite induction
keeps all activity arguments admissible independently
of the scalar value.

Joint real-coupling and quotient continuity, together with
continuity of the maximum in
\eqref{fts-clipped-evaluation}, makes every finite clipped
history continuous in its initial $\gamma$. On a compact
initial interval, \eqref{fts-absolute-increment} bounds all
intermediate scalars. Hence their evaluated couplings
belong to a compact positive interval. Neither continuity
at zero coupling nor a continuous representative selection
nor a volume-uniform modulus is needed.

Since \eqref{fts-decrement-error} and
\eqref{fts-drift-hypothesis} apply to every admissible
realized weak sequence, the clipped history also obeys
\eqref{fts-block-decrements}.

\subsection{Endpoint tuning and removal of clipping}
\label{fts-tuning-section}

\begin{theorem}[Conditional finite ultraviolet tuning]
\label{fts-theorem}
Assume \eqref{fts-drift-hypothesis} with the uniformity just
specified. Fix $m\geq1$ and a permitted finite history of
$mq$ steps whose final period is above the strengthened
threshold, with the nested-family compatibility conditions.
Fix a terminal inverse coupling satisfying
\begin{equation}
       \Gamma_*\geq\Gamma_{\mathrm w}+qD_{\mathrm{abs}}.
 \label{fts-terminal-margin}
\end{equation}
For $x_0=0$, or any fixed real extracted input in the
endpoint ball, some initial inverse coupling in
\begin{equation}
       [\,\Gamma_*+mB_-,\ \Gamma_*+mB_+\,]
 \label{fts-shooting-interval}
\end{equation}
generates $mq$ genuine weak local extraction steps and
ends exactly at $\Gamma_*$. The actual scalar relation
is $\gamma_{j+1}=\gamma_j-\beta_{j+1}(F_{j,g_j}(x_j))$,
with $g_j=\gamma_j^{-1/2}$ throughout.
\end{theorem}

\begin{proof}
Run the clipped recursion with the fixed initial activity.
Summing \eqref{fts-block-decrements}, the final scalar
at the left endpoint of \eqref{fts-shooting-interval}
is at most $\Gamma_*$, while at the right endpoint it
is at least $\Gamma_*$. Continuity and the intermediate
value theorem give a clipped history ending at $\Gamma_*$.

For this selected history, sum the positive \emph{future}
block decrements to obtain
\begin{equation}
 \Gamma_*+(m-k)B_-\leq\gamma_{kq}
                   \leq\Gamma_*+(m-k)B_+,
 \qquad 0\leq k\leq m.
 \label{fts-endpoint-bounds}
\end{equation}
Within any block, \eqref{fts-absolute-increment} implies
\begin{equation}
 \gamma_{kq+r}\geq\gamma_{kq}-rD_{\mathrm{abs}}
       \geq\Gamma_*-qD_{\mathrm{abs}}
       \geq\Gamma_{\mathrm w},
 \qquad 0\leq r\leq q.
 \label{fts-no-clipping}
\end{equation}
Thus clipping never occurred on the selected history.
Its scalar and activity updates are the actual local
extraction updates, and its coupling stays weak.
The exact finite representative and vacuum compensation
identities implement the same normalized local density
at every step.
\end{proof}

The initial bare coupling therefore obeys
\begin{equation}
 (\Gamma_*+mB_+)^{-1/2}
       \leq g_{\mathrm{initial}}
       \leq(\Gamma_*+mB_-)^{-1/2}.
 \label{fts-bare-coupling-bounds}
\end{equation}
Conditional on the drift hypothesis, this tunes arbitrarily
long finite histories to one fixed positive weak endpoint,
with bare coupling of order $m^{-1/2}$. It proves no
uniqueness of the tuned initial value or consistency between
the choices at different $m$.

\subsection{What remains unproved}
\label{fts-scope-section}

The additional input is precisely
\eqref{fts-drift-hypothesis} for the actual finite
determinant coefficients with their actual extraction.
The established activity bounds, parity, continuity and
qualitative finite leakage do not prove it. A determinant
cocycle can telescope the infinite counterpart, but a
quantitative one-loop comparison and finite-regulator
control must still yield the required positive averaged
finite drift.

The horizons considered here are multiples of $q$, enough
for an unbounded ultraviolet subsequence if that drift
input is supplied. No assertion about every integer
horizon is needed. The initial period grows to keep the
final period above the common threshold. The theorem
does not establish convergence of the tuned measures,
mixed or large-field histories, a nontrivial continuum
limit, or reconstruction. Its density statement remains
the one for the specified homogeneous local extraction
map and its justified continuation.


\section{The remaining determinant estimate and the existence argument}
\label{dfn-next}

The leading correction is now an actual functional,
\[
 D_{M,R}=\Delta_{M,R}+\mathcal Q_{\mathrm{det},M,R},
\]
with an exact Hessian trace formula and an exact finite
composition law. In a jointly compatible infinite system its
Maxwell coefficient satisfies
\[
 \beta_{MRL}(D_{M,RL})
   =\beta_{MR}(D_{M,R})+\beta_{MRL}(D_{MR,L}).
\]
These identities supply the algebra for accumulated drift;
they do not determine its value or sign.

The exact gauge and coarea bridge makes the comparison
operators concrete. With its stated finite coordinates,
\[
 D_{M,K}=\mathcal N_0\!\left[
    \log\det G_b-\frac12\log\det Q_b
                 -\frac12\log\det(CQ_b^{-1}C^T)\right],
\]
where \(\mathcal N_0\) subtracts the value at the identity.
The inverse pivot has been accounted for by coarea.
In fact \(CQ_b^{-1}C^T=(q_{\mathrm{eff}}+B)^{-1}\).
Choosing one fixed sufficiently large \(B=aI\), the
coarse term \(\tfrac12\mathcal N_0\log\det(q_{\mathrm{eff}}+aI)\)
has a uniformly bounded Maxwell coefficient by the
proved coarse-action support bounds.
The remaining comparison concerns the fine combination
\[
 \mathcal N_0\!\left[
 \log\det G_b-\frac12\log\det
 \bigl(s''-\lambda\cdot f''+D_bD_b^T+aC^TC\bigr)\right].
\]
In particular, the point-anchor Rayleigh test rules out a
uniform physical mass estimate for the based ghost alone.

A concrete route is to compare the composite determinant from
the bare Wilson action to a background-field vector and ghost
reference at the same cutoff and infrared scale. With every
normalization specified, the desired estimate has the form
\[
 \bigl|\beta(D_{1,R^k})-B_k^{\mathrm{ref}}\bigr|\le C,
 \qquad
 \bigl|B_k^{\mathrm{ref}}-a\,k\bigr|\le C',
 \qquad a>0,
\]
with constants independent of \(k\). Neither displayed bound is
proved in this notebook. The stationary lift, constrained-current
curvature and measure factors must be included in that comparison.
A coefficient quoted for a different effective action does not
establish these estimates for the present one.

The drift hypothesis required for shooting concerns the actual
finite regulators, including arbitrary permitted preceding depths.
An infinite comparison must therefore retain the uniformity and
joint-limit compatibility needed to transfer a positive block
margin to those finite coefficients. The qualitative transfer
proposition supplies that passage once a positive margin holds
in every relevant structural limit; a quantitative
finite-volume convergence rate is unnecessary for this purpose.
No exact finite commutation of Maxwell extraction with classical
transport has been assumed.

Once that drift is supplied, the conditional shooting theorem
closes the finite weak-coupling trajectory issue for this local
map. The subsequent existence argument must still incorporate
all required field histories with their exact measure identities,
prove uniform source estimates, construct a nontrivial continuum
observable hierarchy, and establish its Euclidean symmetry,
positivity and reconstruction properties. These requirements are
those of Notebook 01; the determinant and tuning arguments do not
replace them.

\begin{thebibliography}{9}
\bibitem{priorflow}
\emph{Toward a four-dimensional Yang--Mills construction},
research Notebooks 15--31 in this workspace.
The local minimum, exact constrained integral, common-domain
quantum map, finite induction and classical Maxwell transport
are the principal inputs.
\bibitem{lwbackground}
M.~L\"uscher and P.~Weisz,
\emph{Background field technique and renormalization in lattice
gauge theory}, Nucl. Phys. B \textbf{452} (1995), 213--233.
\href{https://arxiv.org/abs/hep-lat/9504006}{Primary preprint}.
\bibitem{lwmatching}
M.~L\"uscher and P.~Weisz,
\emph{Computation of the relation between the bare lattice
coupling and the MS coupling in SU(N) gauge theories to two
loops}, Nucl. Phys. B \textbf{452} (1995), 234--260.
\href{https://arxiv.org/abs/hep-lat/9505011}{Primary preprint}.
These papers supply background-field formalism and reference
amplitudes to investigate. They are not used here as a proof of
the required comparison with our constrained determinant.
\bibitem{sourceflow}
E.~Sussman and collaborators,
\href{https://github.com/ethanwsussman/Yang--Mills-LLM-experiment}
{Yang--Mills-LLM-experiment}, deposited raw accounts at commit
\texttt{e392ef4501\allowbreak 3a581664ba\allowbreak b76b908671\allowbreak 70b30fe10e}.
\end{thebibliography}
\end{document}
